% concatenated from unfol_derived_fol.tex
\newcommand\hmmax{0}
\newcommand\bmmax{0}
\pdfoutput=1

\newif\ifpersonal
\personalfalse  

\documentclass[11pt,reqno,a4paper,dvipsnames,x11names]{amsart}

% ------------------------------------------------------------
% Encoding / language / layout
% ------------------------------------------------------------
\usepackage[utf8]{inputenc}
\usepackage[T1]{fontenc}
\usepackage[english]{babel}
\usepackage[a4paper,margin=1.12in]{geometry}

% ------------------------------------------------------------
% Fonts / microtype
% ------------------------------------------------------------
\usepackage{microtype}
\usepackage{mathpazo}
\usepackage{euler}

% ------------------------------------------------------------
% Math
% ------------------------------------------------------------
\usepackage{amsmath,amssymb,amsthm,mathtools,mathrsfs,bm,eucal,tensor}
\usepackage{stmaryrd}
\usepackage{dsfont}

% ------------------------------------------------------------
% Diagrams / TikZ
% ------------------------------------------------------------
\usepackage[all,cmtip]{xy}
\usepackage{tikz}
\usetikzlibrary{
 arrows.meta,
 calc,
 positioning,
 fit,
 backgrounds,
 shapes.geometric,
 decorations.markings
}
\usepackage{tikz-cd}

\tikzset{
 >=Stealth,
 box/.style={draw, rounded corners, inner sep=6pt},
 v/.style={draw, circle, inner sep=1.4pt},
 arr/.style={->, thick},
 darr/.style={->, thick, dashed},
 midarrow/.style={
 postaction={decorate},
 decoration={markings,mark=at position 0.58 with {\arrow{Stealth}}}
 }
}

% ------------------------------------------------------------
% Utilities
% ------------------------------------------------------------
\usepackage{enumitem}
\usepackage{comment,braket,xspace,csquotes}

% ------------------------------------------------------------
% Hyperlinks / references
% ------------------------------------------------------------
\usepackage{xcolor}
\usepackage{hyperref}
\hypersetup{
 colorlinks=true,
 linkcolor=blue!60!black,
 citecolor=blue!60!black,
 urlcolor=blue!60!black
}
\usepackage[capitalize,nameinlink]{cleveref}

\numberwithin{equation}{section}

% ------------------------------------------------------------
% Theorem environments
% ------------------------------------------------------------
\theoremstyle{plain}
\newtheorem{theorem}{Theorem}[section]
\newtheorem{proposition}[theorem]{Proposition}
\newtheorem{corollary}[theorem]{Corollary}
\newtheorem{lemma}[theorem]{Lemma}
\newtheorem*{theoremA}{Theorem A}
\newtheorem*{theoremB}{Theorem B}
\newtheorem*{theoremC}{Theorem C}
\newtheorem*{theoremD}{Theorem D}

\theoremstyle{definition}
\newtheorem{definition}[theorem]{Definition}
\newtheorem{example}[theorem]{Example}
\newtheorem{construction}[theorem]{Construction}
\newtheorem{warning}[theorem]{Warning}
\newtheorem{hypothesis}[theorem]{Hypothesis}

\theoremstyle{remark}
\newtheorem{remark}[theorem]{Remark}

% ------------------------------------------------------------
% Notation (keep minimal; compile-safe)
% ------------------------------------------------------------
\newcommand{\C}{\mathbb{C}}
\newcommand{\R}{\mathbb{R}}
\newcommand{\Z}{\mathbb{Z}}
\newcommand{\N}{\mathbb{N}}
\newcommand{\Q}{\mathbb{Q}}
\newcommand{\PP}{\mathbb{P}}

\newcommand{\cO}{\mathcal{O}}
\newcommand{\cM}{\mathcal{M}}
\newcommand{\fol}{\mathcal{F}}
\newcommand{\cB}{\mathcal{B}}
\newcommand{\Aff}{\operatorname{Aff}}
\newcommand{\im}{\mathrm{im}}
\newcommand{\End}{\mathrm{End}}
\newcommand{\Aut}{\mathrm{Aut}}
\newcommand{\Hom}{\mathrm{Hom}}
\newcommand{\rk}{\mathrm{rk}}
\newcommand{\GL}{\mathrm{GL}}
\newcommand{\PGL}{\mathrm{PGL}}
\newcommand{\Div}{\mathrm{Div}}
\newcommand{\Res}{\mathrm{Res}}
\newcommand{\pr}{\operatorname{pr}}

\newcommand{\Gm}{\C^*}
\newcommand{\Sone}{S^1}

% ------------------------------------------------------------
% Article-specific notation
% ------------------------------------------------------------
\DeclareMathOperator{\Spec}{Spec}
\DeclareMathOperator{\Sym}{Sym}
\DeclareMathOperator{\Der}{Der}
\DeclareMathOperator{\RDer}{RDer}
\DeclareMathOperator{\RHom}{RHom}
\DeclareMathOperator{\cofib}{cofib}
\DeclareMathOperator{\hocofib}{hocofib}
\DeclareMathOperator{\fib}{fib}
\DeclareMathOperator{\Tot}{Tot}
\DeclareMathOperator{\MC}{MC}
\DeclareMathOperator{\CE}{CE}
\DeclareMathOperator{\DR}{DR}
\DeclareMathOperator{\Diff}{Diff}
\DeclareMathOperator{\Fol}{Fol}
\DeclareMathOperator{\Unf}{Unf}
\DeclareMathOperator{\Flat}{Flat}
\DeclareMathOperator{\Map}{Map}
\DeclareMathOperator{\Crys}{Crys}
\DeclareMathOperator{\QCoh}{QCoh}
\DeclareMathOperator{\id}{id}

\newcommand{\kk}{k}
\newcommand{\cA}{\mathcal A}
\newcommand{\cE}{\mathcal E}
\newcommand{\cF}{\mathcal F}
\newcommand{\cG}{\mathcal G}
\newcommand{\cR}{\mathcal R}
\newcommand{\g}{\mathfrak g}
\newcommand{\h}{\mathfrak h}
\newcommand{\e}{\mathfrak e}
\newcommand{\kfr}{\mathfrak k}
\newcommand{\z}{\mathfrak z}
\newcommand{\uuder}{\mathfrak u_{\mathrm{der}}}
\newcommand{\Uder}{\mathbb U_{\mathrm{der}}}
\newcommand{\bbT}{\mathbb T}
\newcommand{\bbL}{\mathbb L}
\newcommand{\eps}{\varepsilon}
\newcommand{\qis}{\simeq}
\newcommand{\xto}[1]{\xrightarrow{#1}}
\newcommand{\dg}{\mathrm{dg}}
\newcommand{\gm}{\mathrm{gm}}
\newcommand{\bas}{\mathrm{bas}}
\newcommand{\der}{\mathrm{der}}
\newcommand{\tr}{\mathrm{tr}}
\newcommand{\str}{\mathrm{str}}
\newcommand{\eff}{\mathrm{eff}}
\newcommand{\Lie}{\mathrm{Lie}}
\newcommand{\from}{\colon}
\newcommand{\holim}{\operatorname*{holim}}
\newcommand{\hocolim}{\operatorname*{hocolim}}
\newcommand{\Tate}{\mathrm{Tate}}
\newcommand{\GM}{\mathrm{GM}}
\title[Transversal unfoldings of derived foliations]{Transversal unfoldings of derived foliations}
\author{}
\date{}

\begin{document}



\begin{abstract}
We establish a homotopy-invariant theory of transversal unfoldings of relative derived foliations. A transversal unfolding is an integrable lift of the parameter directions to infinitesimal transverse symmetries of the relative foliation. We construct a canonical transverse controller and prove that the groupoid of strict transversal unfoldings is the one-truncation of its derived space of flat splittings. For affine Chevalley--Eilenberg presentations, the controller is described by first-order Lie derivations modulo inner derivations. When the inner action has a kernel, the ordinary quotient is replaced by a crossed controller whose derived mapping space retains the transported central isotropy and the ensuing higher homotopies. We then construct the controller intrinsically from the graded-mixed de Rham algebra, prove its independence of the chosen presentation, and obtain a global classification under explicit descent hypotheses. The theory encompasses the classical, logarithmic, shifted-Poisson, and representable leaf-space cases.
\end{abstract}
\author[M. Corr\^ea]{Maur\'icio Corr\^ea}
\address[M. Corr\^ea]{Dipartimento di Matematica, Universit\`a degli Studi di Bari, Via E. Orabona 4, I-70125, Bari, Italy}
\email{mauricio.correa.mat@gmail.com, mauricio.barros@uniba.it}

\maketitle
 

\section{Introduction}

Derived foliations, introduced by To\"en and Vezzosi, encode possibly singular foliations by graded-mixed de Rham algebras and retain homotopical information absent from their classical truncations \cite{TV}. They belong naturally to derived algebraic geometry \cite{HAGII,PTVV,CPTVV}, and give rise to derived foliated de Rham cohomology and formal leaves \cite[Sections~2.3 and~3.2--3.3]{TV}. In the rigid holomorphic setting, singular Frobenius theory further relates them to analytic leaf spaces \cite{Malgrange76,Malgrange77,TVRH}. The purpose of the present paper is to determine the infinitesimal structure governing transverse transport in a family of relative derived foliations.

Let \(\pi\from X\to S\) be smooth and let \(\cF/S\) be a relative derived foliation. A deformation records the variation of the leafwise structure along \(S\). A \emph{transversal unfolding} includes, in addition, a lift of the tangent directions of the base to infinitesimal transverse symmetries of \(\cF/S\), subject to an integrability condition. Equivalently, it defines flat transverse transport along the base. On cotangent complexes, transversality is expressed by a fibre sequence
\[
 \pi^*\bbL_S\longrightarrow\bbL_{\widetilde\cF}
 \longrightarrow\bbL_{\cF/S}
 \longrightarrow\pi^*\bbL_S[1],
\]
where \(\widetilde\cF\) is the absolute derived foliation determined by the unfolding. Dually, the quotient of its tangent complex by the relative tangent complex is \(\pi^*\bbT_S\). This description rests upon the relative cotangent theory and its compatibility with base change \cite{IllusieI,TV}.

The distinction between a deformation and an unfolding is already present in the classical theory. If a codimension-one foliation is locally defined by an integrable form \(\omega\), an infinitesimal relative deformation has the form
$
 \omega_\eps=\omega+\eps\eta,
$
whereas an unfolding contains an additional transverse term
$
 \widetilde\omega_\eps=\omega+\eps\eta+h\,d\eps.
$
In Suwa's theory, the term \(h\,d\eps\) records a lift of the parameter direction to the total space \cite{Suwa81,Suwa83}. The integrability equation asserts that the lifted directions act by transverse infinitesimal symmetries with vanishing curvature. The derived theory developed here is governed by the same geometric principle: the local transverse term is replaced by a flat splitting of a canonical controller. In this respect, the construction extends the classical and Lie-algebroid descriptions of transversal unfoldings in \cite{Quallbrunn,CMQ}.

\medskip
 
Let \(B\to A\) be a smooth morphism between smooth commutative \(\kk\)-algebras, and let
\[
 \rho\from\g\longrightarrow T_{A/B}
\]
be a cofibrant strictly perfect relative dg-Lie algebroid. Its Chevalley--Eilenberg algebra
\[
 \CE^*(\g)=\Sym_A(\g^\vee[1])
\]
presents a relative derived foliation, by the Lie--Rinehart and dg-Lie-algebroid models of the graded-mixed de Rham algebra \cite{Rinehart63,Nuiten,TV}. A \emph{basic first-order derivation} of \(\g\) is an infinitesimal Lie symmetry whose transverse symbol is induced by a vector field on the base; this is the transverse analogue of the Atiyah--Lie--Rinehart algebra of first-order operators \cite{Atiyah57,Rinehart63,CMQ}. These derivations form a dg-Lie algebra \(D^1_{\bas}(\g/B)\), and tangent symmetries define the inner-action morphism
\[
 \iota\from\g\longrightarrow D^1_{\bas}(\g/B),
 \qquad
 x\longmapsto\bigl(\rho(x),\operatorname{ad}_x,0\bigr).
\]
The presentation is called \emph{effective} when \(\iota\) is injective. In this case, set
\[
 \uuder(\g/B)
 :=D^1_{\bas}(\g/B)\big/\iota(\g).
\]
This is a dg-Lie algebroid over \(B\), with anchor
$
 a_{\uuder}\from\uuder(\g/B)\longrightarrow T_{B/\kk}.$
 
A \emph{flat splitting} is a \(B\)-linear section of this anchor which is a morphism of dg-Lie algebroids. It is therefore an integrable choice of lifts of vector fields on the base.

The first main result is the affine classification in the absence of isotropy.

\begin{theoremA} 
Let \(B\to A\) be smooth, and let \(\g\to T_{A/B}\) be a cofibrant strictly perfect effective relative dg-Lie algebroid. The effective unfolding groupoid has no non-trivial stabilisers, and the adjoint-symbol construction induces a bijection
\[
 \left|\Unf^{\tr,\eff}(\g/B)\right|
 \simeq
 \Flat\bigl(T_{B/\kk},\uuder(\g/B)\bigr).
\]
Moreover, the basic part of an unfolding reconstructs the full absolute dg-Lie algebroid after extension of scalars from \(B\) to \(A\); no separate basic-generation or faithful-flatness hypothesis is required.
\end{theoremA}

This statement is proved in Theorem~\ref{thm:strict-affine}, together with Lemmas~\ref{lem:automatic-basic-reconstruction} and~\ref{lem:reconstruction-basic-extension}. Given an unfolding, a basic lift acts on \(\g\) by the adjoint representation and therefore determines a basic first-order derivation. Changing the lift changes this derivation by an inner one, while the Jacobi identity gives the flatness of the resulting splitting. Conversely, a flat splitting determines a basic extension over \(B\), and the Lie--Rinehart formula reconstructs the absolute algebroid over \(A\). In this form, Theorem~A is the Chevalley--Eilenberg counterpart of the classical correspondence between unfoldings and flat transverse infinitesimal symmetries \cite{Suwa81,Suwa83,Quallbrunn,CMQ}.

\medskip
 
Effectivity is not intrinsic. If \(\iota\) has a kernel, then
$
 \ker(\iota)
 =\{x\in\g:\rho(x)=0,\ \operatorname{ad}_x=0\}
$
is the central isotropy. Passage to the ordinary quotient would discard these elements, although their transport along the base contributes to the automorphism complex of an unfolding. The appropriate controller is therefore the crossed module
\[
 \Uder(\g/B)
 :=\bigl[\g\xto{\iota}D^1_{\bas}(\g/B)\bigr].
\]
Crossed modules give the customary strict presentation of the corresponding two-term homotopical quotient \cite{BrownSpencer76,Noohi07}. We define its derived space of flat splittings by
\[
 \operatorname{Flat}^{\mathrm{der}}
 \bigl(T_{B/\kk},\Uder(\g/B)\bigr)
 :=
 \Map_{(\mathrm{LieAlgd}^{\infty}_{B})_{/T_{B/\kk}}}
 \bigl(T_{B/\kk},\Uder(\g/B)\bigr),
\]
where the mapping space is computed in the localised \(\infty\)-category of dg-Lie algebroids, or equivalently of \(L_\infty\)-algebroids, over \(T_{B/\kk}\) \cite{Hinich01,HAGII,Pridham10,Nuiten}. As usual, if \(\mathcal G\) is a groupoid, we denote its simplicial nerve by \(N\mathcal G\). Thus \(N\operatorname{Unf}^{\tr}_{\mathrm{str}}(\g/B)\) is a one-type, whereas the derived mapping space also retains the higher homotopies arising from central isotropy.

\begin{theoremB} 
For a cofibrant strictly perfect CE presentation, the adjoint-symbol construction identifies the strict unfolding groupoid with the strict flat-lift groupoid of the crossed controller. Its nerve is the one-truncation of the derived splitting space:
\[
 N\operatorname{Unf}^{\tr}_{\mathrm{str}}(\g/B)
 \simeq
 \tau_{\leq1}\operatorname{Flat}^{\mathrm{der}}
 \bigl(T_{B/\kk},\Uder(\g/B)\bigr).
\]
The full derived space is the mapping space before truncation. The automorphism group of a strict splitting is the Chevalley--Eilenberg cocycle group with coefficients in its transported central isotropy. When \(\iota\) is injective, Theorem~A is recovered upon passing to connected components.
\end{theoremB}

Thus the passage from effective to non-effective unfoldings is not obtained by discarding isotropy, but by replacing an ordinary quotient with its homotopical quotient.

\medskip
 
The preceding construction is initially expressed in a chosen CE model. Its passage from a strict presentation to an invariant graded-mixed construction follows the homotopical theory of derived foliations and dg-Lie algebroids \cite{HAGII,Nuiten,TV}. For a perfect relative derived foliation, the intrinsic controller is defined by
\[
 \Uder(\cF/S)
 :=
 \bigl[
   \bbT_{\cF/S}
   \longrightarrow
   \Der^{\gm}_{\bas}(\DR(\cF/S))
 \bigr].
\]
Here \(\Der^{\gm}_{\bas}(\DR(\cF/S))\) denotes the basic weight-zero graded-mixed derivations of the relative de Rham algebra; ``basic'' means that the transverse symbol is induced by a vector field on \(S\). The definition depends only upon the graded-mixed de Rham algebra of the foliation.

\begin{theoremC} 
Let \(\cF/S\) be a perfect relative derived foliation admitting a cofibrant strictly perfect CE presentation by \(\g\). Then
\[
 \Der^{\gm}_{\bas}(\CE^*(\g))
 \cong D^1_{\bas}(\g/B),
\]
and the intrinsic inner action corresponds to \(\iota\from\g\to D^1_{\bas}(\g/B)\). Hence the crossed-module controller of the presentation represents \(\Uder(\cF/S)\). Any two such CE presentations determine equivalent two-term \(L_\infty\)-controllers. In the effective case their homotopy cofibres are equivalent, and the ordinary quotients represent them whenever the inner maps are cofibrations. Consequently, the classification by flat splittings is independent of the chosen CE presentation.
\end{theoremC}

The comparison with a CE presentation is Proposition~\ref{prop:intrinsic-comparison}; homotopy invariance and presentation independence are established in Theorem~\ref{thm:presentation-independence} and Corollary~\ref{cor:classification-invariance}. It follows that the controller is not merely an auxiliary construction attached to a strict algebraic model, but an invariant of the derived foliation itself.

\medskip
 
Descent is invoked only after the local classification and the intrinsic controller have been established. Under explicit hypotheses of CE-presentability, perfectness, and descent for the controller, the local splitting spaces glue over a smooth affine hypercover by homotopical descent for derived stacks and derived foliations \cite{HAGII,TV}. The cotangent and tangent constructions are compatible with the usual derived base-change theory \cite{IllusieI}, and the descended controller admits the required derived pushforward to the base.

\begin{theoremD} 
Under Hypothesis~\ref{hyp:global-descent}, transversal derived unfoldings of \(\cF/S\) are equivalent to the derived space of flat splittings of the descended controller:
\[
 \operatorname{Unf}^{\tr}_{\mathrm{der}}(\cF/S)
 \simeq
 \operatorname{Flat}^{\mathrm{der}}_S
 \bigl(\bbT_S,\Uder(\cF/S)\bigr)
 :=
 \Map_{\mathrm{LieAlgdStack}_{/\bbT_S}}
 \bigl(\bbT_S,\mathrm R\pi_*\Uder(\cF/S)\bigr),
\]
with the crossed mapping space in the non-effective case. In the effective strict case, this equivalence induces a bijection on connected components with ordinary flat splittings of the descended effective controller.
\end{theoremD}

This is Theorem~\ref{thm:intrinsic}. Proposition~\ref{prop:checkable-global-descent} records a concrete sufficient criterion for its hypotheses, while Theorem~\ref{thm:descent-controller} proves \'etale descent in the strictly CE-presented scheme-theoretic setting. Thus the global theorem is the globalisation of the controller classification, rather than its definition or point of departure.

\medskip
 
The controller also organises the obstruction and deformation theory of unfoldings. Its anchor sequence carries an Atiyah--Kodaira--Spencer class, in the tradition of Atiyah extensions and cotangent-complex obstruction theory \cite{Atiyah57,IllusieI}; the vanishing of this class is equivalent to the existence of a transverse splitting, and the curvature of a chosen splitting is the obstruction to flatness. Around a flat splitting, infinitesimal deformations are governed by the total Chevalley--Eilenberg dg-Lie algebra of base forms with values in the vertical dg kernel, as in the standard dg-Lie formulation of derived deformation problems \cite{Hinich01,Pridham10,Nuiten}. In a split model, the absolute mixed differential is obtained from the relative mixed differential by adding the base de Rham differential and a transverse Lie-derivative term, and the compatibility with both the internal and mixed differentials is the corresponding total Maurer--Cartan equation.

When a relative leaf space is represented, the abstract splitting becomes a flat Ehresmann connection on the leaf-space fibration, consistently with the representability results for derived foliations \cite{TVRH,TV}. A coherently Cartan-linearised crystal then carries derived Gauss--Manin transport, extending the classical differentiation of relative de Rham cohomology \cite{KatzOda68}; in the weight-zero situation considered here, the induced connection preserves the relevant filtration \cite{Deligne71,Deligne74} and makes the foliated Chern character horizontal, in accordance with the theory of characteristic classes for derived foliations \cite{TVGRR}. These statements are consequences of the classification, rather than additional components of the definition of an unfolding.

\medskip
 
The construction recovers ordinary smooth foliations and Lie algebroids \cite{Rinehart63,Nuiten}, logarithmic tangent algebroids \cite{Kato89}, Suwa's transverse term \cite{Suwa81,Suwa83}, derived pullbacks and representable or formal leaf spaces \cite{TVRH,TV}, and strict Hamiltonian presentations of shifted Poisson structures \cite{PTVV,CPTVV}. In each case the same principle remains visible: transverse transport is governed by a controller of infinitesimal symmetries, and integrability is the flatness of a splitting. When isotropy is present, the crossed controller records the central isotropy, its transport along the base and the higher homotopies which an ordinary quotient would lose.

\medskip
 
Sections~2--5 develop the strict affine theory and prove Theorem~A. Sections~6--9 construct the derived mapping space, identify the strict one-truncation, treat transported central isotropy, construct the projectable intrinsic controller and establish Theorems~B and~C. Sections~10--11 develop Gauss--Manin transport and descent, culminating in Theorem~D. Sections~12--17 study the deformation complex, obstruction classes, flatness, normal forms, representable leaf spaces and characteristic classes. The final section collects the principal geometric examples.


\medskip\paragraph{\bf Acknowledgments.} The author acknowledges the partial support of the Universit\`a degli Studi di Bari. He is grateful to Ariel Molinuevo, Federico Quallbrunn and Simone Noja for many stimulating and illuminating conversations. He is a member of GNSAGA, a national research group of INdAM.



\section{Derived foliations and CE presentations}

This section records the conventions concerning derived foliations, graded-mixed algebras, and CE presentations. The definitions are those of To\"en--Vezzosi \cite{TV}; the use of cotangent complexes is Illusie's convention \cite{IllusieI}; and the Lie--Rinehart language for algebroids follows the classical algebraic model of Rinehart and its dg form as developed by Nuiten \cite{Rinehart63,Nuiten}. The ground field \(\kk\) has characteristic zero. Strict CE computations are made in cofibrant strictly perfect models.

\subsection{Graded mixed de Rham algebras}

Throughout the paper, a derivation of a graded mixed cdga means a weight-zero derivation preserving the grading and commuting with the mixed structure, unless another weight is explicitly indicated. Accordingly, transverse symmetries of a derived foliation preserve the weight grading, whereas derivations of higher weight belong to the internal deformation complex of the graded-mixed algebra.

Following the conventions for graded mixed complexes used in derived foliation theory \cite[Section~1.2]{TV}, a graded mixed complex is a graded complex \(E=\bigoplus_{p\in\mathbb Z}E^{(p)}\), where \(p\) is the weight, endowed with an internal differential \(d\) of bidegree \((1,0)\) and a mixed differential
\[
 \epsilon\from E^{(p)}\longrightarrow E^{(p+1)}[-1]
\]
satisfying \(d^2=\epsilon^2=d\epsilon+\epsilon d=0\). A graded mixed cdga is a commutative algebra object in this symmetric monoidal category.

For a cdga \(A\), the derived de Rham algebra is, with the conventions of \cite[Section~1.4]{TV} and the cotangent-complex theory of \cite{IllusieI},
\[
 \DR(A/\kk)=\Sym_A(\bbL_{A/\kk}[1])
\]
with mixed differential induced by the universal derivation \(A\to\bbL_{A/\kk}\). If \(B\to A\) is a morphism, the relative de Rham algebra is denoted \(\DR(A/B)\). In the smooth underived case, its underlying graded algebra is \(\Omega^\bullet_{A/B}[\bullet]\).

\begin{definition}
Let \(X=\Spec A\). A derived foliation on \(X\), relative to \(\kk\), is a graded mixed cdga \(\DR(\cF)\) in the sense of To\"en--Vezzosi \cite[Definition~2.1.1.1]{TV}, together with a morphism of graded mixed cdga's
\[
 \DR(A/\kk)\longrightarrow \DR(\cF),
\]
such that the underlying graded cdga of \(\DR(\cF)\) is quasi-isomorphic to \(\Sym_A(\bbL_\cF[1])\) for an \(A\)-dg-module \(\bbL_\cF\). The complex \(\bbL_\cF\) is the cotangent complex of \(\cF\).
\end{definition}

We shall use this definition in its homotopical sense: derived foliations form an \(\infty\)-category of graded-mixed cdga's satisfying the customary hypotheses of quasi-freeness, connectivity, and perfectness whenever required. The strict affine arguments are carried out only after the choice of cofibrant strictly perfect CE models. The intrinsic assertions concern the corresponding objects in the localised homotopy category, rather than arbitrary non-cofibrant representatives.

After replacing \(\DR(\cF)\) by a quasi-free graded model, its weight-one part represents the cotangent complex,
\[
 \bbL_\cF\simeq\DR(\cF)^{(1)}[-1].
\]
This identification is well-defined in the homotopy category, not as a statement about an arbitrary strict model. Following \cite[Corollary~1.6.0.4]{TV}, all higher coherences of the mixed structure are encoded by the graded-mixed cdga rather than displayed as separate homotopies. The morphism \(\DR(A/\kk)\to\DR(\cF)\) gives the anchor on cotangent complexes
\[
 \bbL_{A/\kk}\longrightarrow \bbL_\cF.
\]

\begin{definition}
Let \(B\to A\) be a morphism. In the strict affine relative setting used in this paper, a relative derived foliation over \(A/B\) is a graded mixed cdga under the relative de Rham algebra \(\DR(A/B)\), in the sense suggested by the relative de Rham foliation construction of \cite[Section~2.1]{TV}. Thus it is a graded mixed cdga
\[
 \DR(A/B)\longrightarrow \DR(\cF_{A/B})
\]
whose underlying graded cdga is quasi-isomorphic to \(\Sym_A(\bbL_{\cF_{A/B}}[1])\). This is the strict affine relative model used for the unfoldings below; it is not intended to replace the absolute definition of To\"en--Vezzosi on general derived stacks.
\end{definition}

\begin{example}
The identity map
$
 \DR(A/B)\longrightarrow \DR(A/B)$
defines the relative de Rham foliation. Its cotangent complex is \(\bbL_{A/B}\). Geometrically it is the foliation by the fibres of \(\Spec A\to\Spec B\).
\end{example}

\subsection{dg-Lie algebroids}

Let \(T_{A/\kk}=\Der_\kk(A,A)\). We use the Lie--Rinehart model for dg-Lie algebroids, originating in Rinehart's algebraic formulation of Lie algebroids and used in derived Lie algebroid theory \cite{Rinehart63,Nuiten,TVRH,TV}. A dg-Lie algebroid over \(A/\kk\) is an \(A\)-dg-module \(\g\), together with a \(\kk\)-linear dg-Lie bracket, and an anchor
\[
 \rho\from \g\longrightarrow T_{A/\kk}
\]
which is \(A\)-linear, is a morphism of dg-Lie algebras, and satisfies the Lie--Rinehart identity
\[
 [x,ay]=\rho(x)(a)y+(-1)^{|x||a|}a[x,y]
\]
for homogeneous \(a\in A\) and \(x,y\in\g\). A relative dg-Lie algebroid over \(A/B\) is defined similarly with anchor landing in \(T_{A/B}\).

\begin{construction}
Let \(\g\) be a cofibrant strictly perfect relative dg-Lie algebroid over \(A/B\). Its Chevalley--Eilenberg graded-mixed cdga is the standard comparison algebra between dg-Lie algebroids and derived foliations \cite[Section~2.2]{TV}; see also \cite{Rinehart63,Nuiten} for the Lie--Rinehart and dg-Lie-algebroid background:
\[
 \CE^*(\g):=\Sym_A(\g^\vee[1]).
\]
The Chevalley--Eilenberg mixed differential is determined by the anchor and the bracket, while the internal differential is induced by the differential of \(\g\). In degree one, for \(\alpha\in\g^\vee\) and homogeneous \(x,y\in\g\), the bracket-anchor part of the mixed differential has the form
\[
 \epsilon_{CE}\alpha(x,y)=
 \rho(x)(\alpha(y))
 -(-1)^{|x||y|}\rho(y)(\alpha(x))
 -\alpha([x,y]),
\]
with the standard Koszul signs in the full formula. The dual of the anchor induces
\[
 \DR(A/B)\longrightarrow \CE^*(\g),
\]
so \(\CE^*(\g)\) is a relative derived foliation.
\end{construction}

\begin{definition}
A relative derived foliation \(\cF_{A/B}\) is called strictly CE-presented if there is a cofibrant strictly perfect relative dg-Lie algebroid \(\g\) and an equivalence of graded mixed cdga's under \(\DR(A/B)\)
\[
 \DR(\cF_{A/B})\simeq\CE^*(\g).
\]
In the strict theorem the presentation \(\g\) is included among the hypotheses. An arbitrary perfect CE presentation is first replaced locally by a cofibrant strictly perfect model before the strict affine theorem is applied.
\end{definition}

\begin{warning}
The strict controller associated with a chosen dg-Lie algebroid presentation is not claimed to be independent of that presentation as an ordinary dg-Lie algebroid. The intrinsic controller is the homotopy quotient of the graded-mixed derivations of \(\DR(\cF_{A/B})\), constructed in Section~\ref{sec:intrinsic}. Accordingly, the strict affine classification is formulated in a chosen cofibrant CE model; invariance is obtained only after passage to the graded-mixed controller and the corresponding homotopy quotient.
\end{warning}

\section{Strict transversal unfoldings}

Let \(B\to A\) be a smooth morphism between smooth commutative \(\kk\)-algebras. Write
\[
 q\from T_{A/\kk}\longrightarrow A\otimes_B T_{B/\kk}
\]
for the quotient map in
\[
 0\longrightarrow T_{A/B}\longrightarrow T_{A/\kk}
 \xto{q} A\otimes_B T_{B/\kk}\longrightarrow 0.
\]
There is always the map \(T_{B/\kk}\to A\otimes_B T_{B/\kk}\), \(\xi\mapsto1\otimes\xi\), but it need not be injective unless \(B\to A\) is faithfully flat. We therefore do not regard \(T_{B/\kk}\) as a literal submodule of \(A\otimes_B T_{B/\kk}\). Basic transverse sections are therefore represented by pairs over this map, equivalently by fibre products over \(A\otimes_B T_{B/\kk}\).

\begin{definition}
Let \(\rho\from\g\to T_{A/B}\) be a relative dg-Lie algebroid. A strict transversal unfolding of \(\g\), in the spirit of Suwa's unfolding theory and its Lie-algebroid extension \cite{Suwa81,Suwa83,CMQ}, is the strict algebraic form of an integrable transverse deformation: it is an absolute dg-Lie algebroid
\[
 \widetilde\rho\from\widetilde\g\longrightarrow T_{A/\kk}
\]
together with an inclusion of dg-Lie algebroids \(\g\hookrightarrow\widetilde\g\), such that:
\begin{enumerate}[label=(\roman*)]
\item \(\g\) is a dg-Lie ideal in \(\widetilde\g\);
\item \(\widetilde\rho|_\g=\rho\), viewed as a map to \(T_{A/B}\subset T_{A/\kk}\);
\item there is an isomorphism of \(A\)-dg-modules
$
 \widetilde\g/\g\simeq A\otimes_B T_{B/\kk};$
\item under this isomorphism the induced anchor on the quotient is the identity of \(A\otimes_B T_{B/\kk}\).
\end{enumerate}
\end{definition}

Throughout the paper, morphisms of strict transversal unfoldings are understood in this rigid sense: they are isomorphisms of absolute dg-Lie algebroids over \(A\) which restrict to the identity on the fixed relative algebroid \(\g\) and induce the identity on the prescribed quotient \(A\otimes_B T_{B/\kk}\). Thus the groupoid records possible stabilisers of an unfolding, but it does not quotient by automorphisms of the original relative foliation. This convention is required in the non-effective case, where central isotropy survives through automorphisms of the unfolding.
The defining diagram is
\[
\begin{tikzcd}[column sep=large,row sep=large]
0 \arrow[r] & \g \arrow[r] \arrow[d,"\rho"'] & \widetilde\g \arrow[r] \arrow[d,"\widetilde\rho"] & A\otimes_B T_{B/\kk} \arrow[r] \arrow[d,equal] & 0 \\
0 \arrow[r] & T_{A/B} \arrow[r] & T_{A/\kk} \arrow[r,"q"] & A\otimes_B T_{B/\kk} \arrow[r] & 0 .
\end{tikzcd}
\]

\begin{definition}
Let \(\cF_{A/B}\) be a strictly CE-presented relative derived foliation represented by \(\g\). A strict CE-presented transversal unfolding of \((\cF_{A/B},\g)\) is a strict transversal unfolding \(\widetilde\g\) of \(\g\). The associated absolute derived foliation is
\[
 \DR(\widetilde\cF):=\CE^*(\widetilde\g).
\]
Its cotangent complex fits into the triangle
\[
 A\otimes_B\bbL_{B/\kk}
 \longrightarrow
 \bbL_{\widetilde\cF}
 \longrightarrow
 \bbL_{\cF_{A/B}}
 \longrightarrow
 (A\otimes_B\bbL_{B/\kk})[1].
\]
\end{definition}

\begin{remark}
The cotangent triangle is the derived replacement of the rank condition. It says that the transverse directions of the absolute foliation are the directions of the base.
\end{remark}

\section{Basic first-order derivations}

The first-order operator algebra used here is the Lie--Rinehart analogue of Atiyah's algebroid of first-order differential operators \cite{Atiyah57,Rinehart63}. In the present context the operators must also preserve the dg-Lie bracket and have symbols compatible with the anchor, so that they act on the CE algebra by graded-mixed derivations.

\subsection{The derivation dg-Lie algebra}

\begin{definition}
Let \(\rho\from\g\to T_{A/B}\) be a relative dg-Lie algebroid. The dg-Lie algebra \(D^1_{A/\kk}(\g)\), which is the derived analogue of the first-order operator algebroid used in \cite{Atiyah57,Rinehart63,CMQ}, consists, in degree \(n\), of pairs
\[
 (\theta,\delta),
 \qquad
 \theta\in T_{A/\kk}^n,
 \quad
 \delta\from\g\to\g[n],
\]
satisfying, for homogeneous \(a\in A\) and \(x,y\in\g\),
\[
 \delta(ax)=\theta(a)x+(-1)^{n|a|}a\delta(x),
\]
\[
 \delta([x,y])=[\delta x,y]+(-1)^{n|x|}[x,\delta y],
\]
\[
 \rho(\delta x)=[\theta,\rho(x)].
\]
The differential is the commutator with the internal differentials,
\[
 d(\theta,\delta)=([d_A,\theta],[d_\g,\delta]),
\]
which in the underived algebra case reduces to the internal differential on the endomorphism complex of \(\g\). The bracket is
\[
 [(\theta,\delta),(\theta',\delta')]
 =([\theta,\theta'],[\delta,\delta']).
\]
The symbol is the dg-Lie morphism
\[
 \sigma\from D^1_{A/\kk}(\g)\longrightarrow T_{A/\kk},
 \qquad
 (\theta,\delta)\longmapsto\theta.
\]
\end{definition}

\begin{lemma}
The complex \(D^1_{A/\kk}(\g)\) is a dg-Lie algebra and \(\sigma\) is a morphism of dg-Lie algebras.
\end{lemma}

\begin{proof}
The commutator of two first-order operators is first-order, as in the usual Atiyah--Lie--Rinehart first-order calculus \cite{Atiyah57,Rinehart63}, and its symbol is the commutator of the symbols. The Leibniz identity follows by expanding \([\delta,\delta'](ax)\) and observing that the second-order terms cancel with the Koszul signs. The bracket-derivation identity follows because the commutator of two derivations of a dg-Lie algebra is again a derivation. Finally,
\[
 \rho([\delta,\delta']x)
 =[\theta,[\theta',\rho(x)]]
 -(-1)^{|\theta||\theta'|}[\theta',[\theta,\rho(x)]],
\]
which is \([[\theta,\theta'],\rho(x)]\) by the graded Jacobi identity. Hence anchor-compatibility is preserved. Compatibility with the internal differential is the same calculation with one of the two derivations equal to the differential.
\end{proof}

\begin{definition}
The inner derivation map is
$
 \iota\from \g\longrightarrow D^1_{A/\kk}(\g),
 $ given by $
 x\longmapsto (\rho(x),\operatorname{ad}_x).
$
The relative dg-Lie algebroid \(\g\) is called effective if \(\iota\) is a monomorphism of complexes.
\end{definition}

\begin{warning}
Effectivity is a property of the chosen strict model. Intrinsically, the transverse controller is the crossed module, or homotopy quotient; the ordinary effective quotient is used only when the inner-action morphism is represented by a monomorphism in the chosen cofibrant model.
\end{warning}

\begin{lemma}
The map \(\iota\) is a morphism of dg-Lie algebras. Its kernel is the complex
\[
 \z(\g):=\{x\in\g: \rho(x)=0,\ \operatorname{ad}_x=0\}.
\]
Moreover, the image of \(\iota\) is a dg-Lie ideal in every dg-Lie subalgebra of \(D^1_{A/\kk}(\g)\) stable under the action on \(\g\).
\end{lemma}

\begin{proof}
The Lie--Rinehart identity says that \(\operatorname{ad}_x\) is a first-order operator with symbol \(\rho(x)\). The equality
\[
 [\operatorname{ad}_x,\operatorname{ad}_y]=\operatorname{ad}_{[x,y]}
\]
is the graded Jacobi identity, and \(\rho([x,y])=[\rho(x),\rho(y)]\). The stated description of the kernel follows directly from the definition. If \((\theta,\delta)\in D^1_{A/\kk}(\g)\), then
\[
 [(\theta,\delta),(\rho(x),\operatorname{ad}_x)]
 =(\rho(\delta x),\operatorname{ad}_{\delta x}),
\]
again by the derivation property and the anchor-compatibility condition. Thus the image is an ideal.
\end{proof}

\subsection{Basic derivations}

\begin{definition}
The dg-Lie algebra of basic first-order derivations of \(\g\) over \(B\) is the sub-dg-Lie algebra
\[
 D^1_{\bas}(\g/B)\subset D^1_{A/\kk}(\g)\times T_{B/\kk}
\]
whose homogeneous sections are triples \((\theta,\delta,\xi)\), with \((\theta,\delta)\in D^1_{A/\kk}(\g)\) and \(\xi\in T_{B/\kk}\), such that
\[
 q(\theta)=1\otimes\xi.
\]
The bracket is induced by
$
 [(\theta,\delta,\xi),(\theta',\delta',\eta)]
 =([\theta,\theta'],[\delta,\delta'],[\xi,\eta]).
$
\end{definition}

Since \(\rho(\g)\subset T_{A/B}\), the inner derivation map lands in the kernel of the projection to \(T_{B/\kk}\). Explicitly, in the basic derivation algebra it is
\[
 \iota(x)=\bigl(\rho(x),\operatorname{ad}_x,0\bigr).
\]
Thus
$
 \iota\from\g\longrightarrow D^1_{\bas}(\g/B)
$
defines a morphism of dg-Lie algebras.

\begin{definition}
Assume \(\g\) is effective. The strict derived transverse symmetry algebroid of \(\g\) over \(B\) is the quotient dg-Lie algebra
\[
 \uuder(\g/B):=D^1_{\bas}(\g/B)/\iota(\g).
\]
It has an anchor
$
 a_{\uuder}\from \uuder(\g/B)\longrightarrow T_{B/\kk}
$
induced by \((\theta,\delta,\xi)\mapsto\xi\).
\end{definition}

\begin{lemma}\label{lem:base-LR-controller}
The \(B\)-module structure on \(D^1_{\bas}(\g/B)\), obtained by restricting scalars along \(B\to A\), makes
\[
 D^1_{\bas}(\g/B)\longrightarrow T_{B/\kk},
 \qquad (\theta,\delta,\xi)\longmapsto \xi,
\]
a dg-Lie algebroid over \(B\). The image of
\(\iota\from\g\to D^1_{\bas}(\g/B)\) is a vertical dg-Lie ideal. Consequently, in the effective case, \(\uuder(\g/B)\) inherits a dg-Lie algebroid structure over \(B\), with anchor \(a_{\uuder}\).
\end{lemma}

\begin{proof}
For \(b\in B\) and \(u=(\theta,\delta,\xi)\), multiplication by \(b\) means multiplication by its image in \(A\) on the first two components and ordinary multiplication on the vector field \(\xi\). If \(v=(\theta',\delta',\eta)\), the commutator calculation in \(D^1_{A/\kk}(\g)\) gives
\[
 [u,bv]=\xi(b)v+b[u,v],
\]
with the usual Koszul sign in the dg case. This is the Lie--Rinehart identity over \(B\), and the anchor is compatible with brackets. Since inner derivations have base component zero, the image of \(\iota\) lies in the kernel of the base anchor. The ideal property was proved above, and quotienting by this vertical ideal gives the asserted structure on \(\uuder(\g/B)\).
\end{proof}

Let \(T^{\bas}_{A/\kk}\subset T_{A/\kk}\times T_{B/\kk}\) denote the Lie algebra of pairs \((Y,\xi)\) such that \(q(Y)=1\otimes\xi\). The construction is summarised by
\[
\begin{tikzcd}[column sep=large,row sep=large]
0 \arrow[r] & \g \arrow[r,"\iota"] \arrow[d,"\rho"'] & D^1_{\bas}(\g/B) \arrow[r] \arrow[d,"\sigma"] & \uuder(\g/B) \arrow[r] \arrow[d,"a_{\uuder}"] & 0 \\
0 \arrow[r] & T_{A/B} \arrow[r] & T^{\bas}_{A/\kk} \arrow[r] & T_{B/\kk} \arrow[r] & 0 .
\end{tikzcd}
\]

\section{The affine strict theorem}

We begin with the classification in the strict affine case. The proof is given explicitly on sections, for it is here that the transverse term of an unfolding appears as an integrability condition rather than an arbitrary deformation parameter. The argument follows the classical pattern of Suwa and its Lie-algebroid reformulation, but with the CE dg-Lie algebroid replacing the ordinary foliation \cite{Suwa81,Suwa83,Quallbrunn,CMQ}.

\begin{definition}
A flat splitting of the anchor
$
 a_{\uuder}\from \uuder(\g/B)\longrightarrow T_{B/\kk}
$
is a \(B\)-linear morphism of dg-Lie algebroids over \(B\)
\[
 s\from T_{B/\kk}\longrightarrow \uuder(\g/B)
\]
whose composite with the anchor is the identity. Equivalently, it is a Lie-algebraic section of the anchor compatible with the \(B\)-module structures. The set of flat splittings is denoted
\[
 \Flat(T_{B/\kk},\uuder(\g/B)).
\]
\end{definition}

\begin{definition}\label{def:basic-part}
Assume that \(B\to A\) is a smooth morphism between smooth commutative \(\kk\)-algebras and that
$
 \rho\from\g\longrightarrow T_{A/B} $
is a cofibrant strictly perfect relative dg-Lie algebroid. Let
\[
 0\longrightarrow \g\longrightarrow\widetilde\g
 \longrightarrow A\otimes_B T_{B/\kk}\longrightarrow0
\]
be a strict transversal unfolding. Its \emph{basic part} is the fibre product
\[
 \widetilde\g^{\bas}
 :=
 \widetilde\g\times_{A\otimes_B T_{B/\kk}}T_{B/\kk},
\]
where \(T_{B/\kk}\to A\otimes_B T_{B/\kk}\) is \(\xi\mapsto1\otimes\xi\). Thus an element of \(\widetilde\g^{\bas}\) is a pair \((\widetilde\xi,\xi)\) such that the image of \(\widetilde\xi\) in \(A\otimes_B T_{B/\kk}\) is \(1\otimes\xi\). The bracket is
\[
 [(\widetilde\xi,\xi),(\widetilde\eta,\eta)]
 :=([\widetilde\xi,\widetilde\eta],[\xi,\eta]),
\]
and the anchor is the pair \((\widetilde\rho(\widetilde\xi),\xi)\). Hence \(\widetilde\g^{\bas}\) is a \(B\)-dg-Lie algebroid fitting into an exact sequence
\[
 0\longrightarrow \g\longrightarrow \widetilde\g^{\bas}
 \longrightarrow T_{B/\kk}\longrightarrow0
\]
where \(\g\) is regarded by restriction of scalars along \(B\to A\). This definition does not require \(T_{B/\kk}\to A\otimes_B T_{B/\kk}\) to be injective. We denote by
\(\Unf^{\tr,\eff}(\g/B)\) the groupoid of strict transversal unfoldings when \(\g\) is effective.
\end{definition}

\begin{lemma}\label{lem:effective-discrete}
Assume that \(\g\) is effective. Then every automorphism of a strict transversal unfolding which is the identity on \(\g\) and on the quotient \(A\otimes_B T_{B/\kk}\) is the identity. Consequently the effective unfolding groupoid is equivalent to a set.
\end{lemma}

\begin{proof}
Let \(\varphi\) be such an automorphism of \(\widetilde\g\). If \((\widetilde\xi,\xi)\in\widetilde\g^{\bas}\), then
\[
 c_\xi:=\varphi(\widetilde\xi)-\widetilde\xi\in\g .
\]
Since \(\varphi\) preserves the anchor and induces the identity on the quotient, \(\rho(c_\xi)=0\). Since \(\varphi\) is the identity on \(\g\), for every \(x\in\g\) one has
\[
 [\widetilde\xi+c_\xi,x]
 =[\varphi(\widetilde\xi),\varphi(x)]
 =\varphi([\widetilde\xi,x])
 =[\widetilde\xi,x].
\]
Thus \(\operatorname{ad}_{c_\xi}=0\). Hence \(c_\xi\in\ker\iota\), and effectivity gives \(c_\xi=0\). Therefore \(\varphi\) is the identity on the basic part. Lemma~\ref{lem:automatic-basic-reconstruction} then implies that \(\varphi\) is the identity on all of \(\widetilde\g\), because the basic part generates \(\widetilde\g\) after extension of scalars from \(B\) to \(A\).
\end{proof}

\begin{lemma}\label{lem:automatic-basic-reconstruction}
Let \(B\to A\) be smooth, and let \(\widetilde\g\) be a strict transversal unfolding of \(\g\). Write \(p_1\from\widetilde\g^{\bas}\to\widetilde\g\) for the projection from the fibre product. Then the map
\[
 \Phi_{\widetilde\g}\from
 (A\otimes_B\widetilde\g^{\bas})/\widetilde{\cR}
 \longrightarrow \widetilde\g,
 \qquad a\otimes u\longmapsto a p_1(u),
\]
where \(\widetilde{\cR}\) is generated by
\[
 a\otimes x-1\otimes ax,
 \qquad a\in A,\quad x\in\g,
\]
is an isomorphism of \(A\)-dg-Lie algebroids. In particular no separate basic-generation hypothesis is needed in the smooth affine classification theorem.
\end{lemma}

\begin{proof}
Surjectivity follows at once. If \(y\in\widetilde\g\), its image in \(A\otimes_BT_{B/\kk}\) may be written as a finite sum \(\sum_i a_i(1\otimes\xi_i)\). Choose pairs \(\widehat\xi_i=(\widetilde\xi_i,\xi_i)\in\widetilde\g^{\bas}\). Then
\[
 y-\sum_i a_i\widetilde\xi_i\in\g,
\]
so \(y\) lies in the image of \(\Phi_{\widetilde\g}\).
For injectivity, let \(\sum_i a_i\otimes u_i\in A\otimes_B\widetilde\g^{\bas}\) map to zero in \(\widetilde\g\). Its image in \(A\otimes_BT_{B/\kk}\) is zero. Since \(B\to A\) is smooth, hence flat, tensoring the exact sequence
\[
 0\longrightarrow\g\longrightarrow\widetilde\g^{\bas}
 \longrightarrow T_{B/\kk}\longrightarrow0
\]
with \(A\) remains exact. Therefore \(\sum_i a_i\otimes u_i\) is congruent, in \(A\otimes_B\widetilde\g^{\bas}\), to an element of \(A\otimes_B\g\). Its image in \(\widetilde\g\) is then the multiplication image in the already \(A\)-linear module \(\g\). The relations defining \(\widetilde{\cR}\) identify \(A\otimes_B\g\) with \(\g\); equivalently they quotient by the kernel of the multiplication map
\[
 A\otimes_B\g\longrightarrow\g,
 \qquad a\otimes x\longmapsto ax .
\]
Hence the original element is zero modulo \(\widetilde{\cR}\). This proves that \(\Phi_{\widetilde\g}\) is an isomorphism of complexes and of \(A\)-modules.

It remains to verify the Lie-algebroid structure. The bracket and anchor on the source are the ones obtained from the fibre-product basic extension by Lemma~\ref{lem:reconstruction-basic-extension}. Under \(\Phi_{\widetilde\g}\), the element \(a\otimes(\widetilde u,u_B)\) is sent to \(a\widetilde u\); the Lie--Rinehart formula in Lemma~\ref{lem:reconstruction-basic-extension} is the restriction of the bracket and anchor of \(\widetilde\g\) to such elements. Hence \(\Phi_{\widetilde\g}\) preserves brackets, anchors, and differentials. This proves the assertion.
\end{proof}


\begin{lemma}\label{lem:reconstruction-basic-extension}
Let \(\e\) be a \(B\)-dg-Lie algebroid fitting into an exact sequence
\[
 0\longrightarrow \g\longrightarrow\e\longrightarrow T_{B/\kk}\longrightarrow0
\]
and suppose that it is equipped with a morphism
$
 \Theta\from \e\longrightarrow D^1_{\bas}(\g/B)
$
whose restriction to \(\g\) is the inner map \(\iota\) and whose projection to \(T_{B/\kk}\) is the quotient map \(\e\to T_{B/\kk}\). Put
\[
 \sigma_\e:=\sigma\circ\Theta\from \e\longrightarrow T_{A/\kk}.
\]
Then
\[
 \widetilde\g_\e:=(A\otimes_B\e)/\cR,
 \qquad
 \cR=\langle a\otimes x-1\otimes ax\mid a\in A,\quad x\in\g\rangle,
\]
has a structure of dg-Lie algebroid over \(A\), with anchor
$
 \widetilde\rho_\e(a\otimes u)=a\sigma_\e(u),
$
and bracket
\[
 [a\otimes u,b\otimes v]
 =ab\otimes[u,v]
 +a\sigma_\e(u)(b)\otimes v
 -(-1)^{|u||v|}b\sigma_\e(v)(a)\otimes u .
\]
Moreover there is an exact sequence
\[
 0\longrightarrow\g\longrightarrow\widetilde\g_\e
 \longrightarrow A\otimes_B T_{B/\kk}\longrightarrow0,
\]
so \(\widetilde\g_\e\) is a strict transversal unfolding.
\end{lemma}

\begin{proof}
The anchor is well-defined because the relation \(a\otimes x=1\otimes ax\) is sent to
\[
 a\rho(x)-\rho(ax)=0,
\]
using the \(A\)-linearity of the relative anchor. The bracket is compatible with the same relation because \(\Theta(u)\) is a first-order Lie derivation of \(\g\). Indeed, for \(x\in\g\) and homogeneous \(b\in A\), the equality
\[
 \Theta(u)(bx)=\sigma_\e(u)(b)x+(-1)^{|u||b|}b\Theta(u)(x)
\]
is the statement that replacing \(b\otimes x\) by \(1\otimes bx\) gives the same class modulo \(\cR\). Compatibility with the relation in the other argument follows from graded skew-symmetry. Hence the displayed formula descends to \(\widetilde\g_\e\).
The Leibniz identity over \(A\) follows from the formula:
\[
 [\alpha, f\beta]
 =\widetilde\rho_\e(\alpha)(f)\beta
 +(-1)^{|\alpha||f|}f[\alpha,\beta]
\]
for local homogeneous sections. The anchor is a morphism of brackets because \(\sigma\from D^1_{\bas}(\g/B)\to T_{A/\kk}\) is a dg-Lie morphism and \(\Theta\) is bracket-preserving. The Jacobiator of the bracket on \(A\otimes_B\e\) is the usual semidirect-product Jacobiator. Its terms split into: the Jacobiator in \(\e\), which vanishes; the Jacobiator of vector fields \(\sigma_\e(u)\), which vanishes because \(\sigma_\e\) is a Lie morphism; and the mixed terms, which cancel exactly because every \(\Theta(u)\) acts on \(\g\) by a derivation of the Lie bracket. Thus \(\widetilde\g_\e\) is a dg-Lie algebroid over \(A\).
The quotient by the embedded copy of \(\g\) is obtained by tensoring the quotient \(\e/\g\simeq T_{B/\kk}\) with \(A\), so it is \(A\otimes_BT_{B/\kk}\). The construction is compatible with the automatic reconstruction isomorphism of Lemma~\ref{lem:automatic-basic-reconstruction}.
\end{proof}

\begin{theorem}\label{thm:strict-affine}
Let \(B\to A\) be a smooth morphism between smooth commutative \(\kk\)-algebras, and let
\[
 \rho\from\g\longrightarrow T_{A/B}
\]
be a cofibrant strictly perfect effective relative dg-Lie algebroid. Then the effective unfolding groupoid has no non-trivial stabilisers, and the adjoint-symbol construction induces a bijection
\[
 \left|\Unf^{\tr,\eff}(\g/B)\right|
 \simeq
 \Flat(T_{B/\kk},\uuder(\g/B)).
\]
\end{theorem}

\begin{proof}
The proof consists of four steps. By Lemma~\ref{lem:effective-discrete}, in the effective case no non-trivial stabiliser remains in the rigid unfolding groupoid, so it is enough to identify isomorphism classes. The construction is the derived analogue of the adjoint-symbol construction for unfoldings of singular Lie algebroids \cite{CMQ}.

\smallskip
\noindent\emph{Step 1: the adjoint-symbol map associated with an unfolding.}
Let
\[
 0\longrightarrow\g\longrightarrow\widetilde\g
 \longrightarrow A\otimes_B T_{B/\kk}\longrightarrow 0
\]
be a strict effective transversal unfolding. Denote by \(\widetilde\g^{\bas}\) the fibre-product basic part of Definition~\ref{def:basic-part}. Thus there is an exact sequence of \(B\)-dg-Lie algebroids
\[
 0\longrightarrow\g\longrightarrow\widetilde\g^{\bas}
 \longrightarrow T_{B/\kk}\longrightarrow0 .
\]
For \(\widehat\xi=(\widetilde\xi,\xi)\in\widetilde\g^{\bas}\), define
\[
 \Theta_{\widehat\xi}:=\bigl(\widetilde\rho(\widetilde\xi),
 \operatorname{ad}_{\widetilde\xi}|_\g\bigr).
\]
We verify that this is an element of \(D^1_{\bas}(\g/B)\). Since \(\g\) is a dg-Lie ideal in \(\widetilde\g\), the operator \(\operatorname{ad}_{\widetilde\xi}|_\g\) preserves \(\g\). For homogeneous \(a\in A\) and \(x\in\g\), the Lie--Rinehart identity in \(\widetilde\g\) gives
\[
 [\widetilde\xi,ax]
 =\widetilde\rho(\widetilde\xi)(a)x
 +(-1)^{|\widetilde\xi||a|}a[\widetilde\xi,x].
\]
Hence \(\operatorname{ad}_{\widetilde\xi}|_\g\) is a first-order differential operator on \(\g\) with symbol \(\widetilde\rho(\widetilde\xi)\). Moreover, the compatibility of the anchor with the bracket gives
\[
 \rho([\widetilde\xi,x])
 =[\widetilde\rho(\widetilde\xi),\rho(x)]
\]
for every \(x\in\g\); therefore the operator is a Lie derivation in the sense defining \(D^1_{A/\kk}(\g)\). If \(\widetilde\xi\) maps to \(\xi\in T_{B/\kk}\), then
\[
 q(\widetilde\rho(\widetilde\xi))=1\otimes\xi.
\]
Thus the symbol is basic and \(\Theta_{\widehat\xi}\in D^1_{\bas}(\g/B)\).
If \(\widehat\xi=(\widetilde\xi,\xi)\) is replaced by \(\widehat\xi+x=(\widetilde\xi+x,\xi)\) with \(x\in\g\), then
\[
 \Theta_{\widehat\xi+x}-\Theta_{\widehat\xi}
 =\bigl(\rho(x),\operatorname{ad}_x,0\bigr)=\iota(x).
\]
Consequently the class of \(\Theta_{\widehat\xi}\) in \(\uuder(\g/B)\) depends only on the image \(\xi\) of \(\widehat\xi\) in \(T_{B/\kk}\). We obtain a well-defined morphism
\[
 s_{\widetilde\g}\from T_{B/\kk}\longrightarrow\uuder(\g/B),
 \qquad
 s_{\widetilde\g}(\xi)=[\Theta_{\widehat\xi}].
\]
By construction its anchor is \(\xi\), so \(a_{\uuder}\circ s_{\widetilde\g}=\id\). The map is \(B\)-linear: if \(b\in B\), then \(\rho(x)(b)=0\) for every relative element \(x\in\g\), and the Lie--Rinehart identity gives
\[
 [b\widetilde\xi,x]=b[\widetilde\xi,x].
\]
Hence \(\Theta_{b\widehat\xi}=b\Theta_{\widehat\xi}\). It is also compatible with the cohomological differentials. Since the quotient \(A\otimes_B T_{B/\kk}\) is concentrated in degree zero, \(d\widetilde\xi\in\g\), and the differential of the first-order derivation \(\Theta_{\widehat\xi}\) is \(\iota(d\widetilde\xi)\). Therefore the class of \(\Theta_{\widehat\xi}\) is closed in the quotient. Thus \(s_{\widetilde\g}\) is a morphism of dg \(B\)-modules over \(T_{B/\kk}\).

\smallskip
\noindent\emph{Step 2: flatness of the splitting.}
Choose basic lifts \(\widehat\xi=(\widetilde\xi,\xi)\) and \(\widehat\eta=(\widetilde\eta,\eta)\) of \(\xi,\eta\in T_{B/\kk}\). Since \(\widetilde\g\to A\otimes_BT_{B/\kk}\) is a morphism of dg-Lie algebras, the element
\[
 [\widetilde\xi,\widetilde\eta]-\widetilde{[\xi,\eta]}
\]
belongs to \(\g\), where \(\widehat{[\xi,\eta]}=(\widetilde{[\xi,\eta]},[\xi,\eta])\) is any basic lift of \([\xi,\eta]\). Applying the adjoint-symbol construction yields
\[
 [\Theta_{\widehat\xi},\Theta_{\widehat\eta}]
 -\Theta_{\widehat{[\xi,\eta]}}
 \in\iota(\g).
\]
After passing to \(\uuder(\g/B)\) this becomes
\[
 [s_{\widetilde\g}(\xi),s_{\widetilde\g}(\eta)]
 =s_{\widetilde\g}([\xi,\eta]).
\]
Thus \(s_{\widetilde\g}\) is a flat splitting.

\smallskip
\noindent\emph{Step 3: reconstruction from a flat splitting.}
Conversely, let
\[
 s\from T_{B/\kk}\longrightarrow\uuder(\g/B)
\]
be a flat \(B\)-linear splitting. Form the pullback dg-Lie algebra
\[
\begin{tikzcd}[column sep=large,row sep=large]
\e_s \arrow[r] \arrow[d] & D^1_{\bas}(\g/B) \arrow[d] \\
T_{B/\kk} \arrow[r,"s"'] & \uuder(\g/B).
\end{tikzcd}
\]
Since \(\iota\) is injective in the effective case, this pullback sits in an exact sequence
\[
 0\longrightarrow\g\longrightarrow\e_s
 \longrightarrow T_{B/\kk}\longrightarrow0.
\]
The composite
\[
 \e_s\longrightarrow D^1_{\bas}(\g/B)
 \xto{\sigma} T_{A/\kk}
\]
will be denoted by \(\sigma_s\). It restricts to \(\rho\) on \(\g\), and its projection to \(A\otimes_BT_{B/\kk}\) is the tautological basic vector field.

It remains to extend \(\e_s\) from \(B\) to \(A\), while identifying the relative ideal \(\g\), which is already \(A\)-linear. Define
\[
 \widetilde\g_s:=(A\otimes_B\e_s)/\cR,
\]
where \(\cR\) is generated by the elements
$
 a\otimes x-1\otimes ax,
 $ with $ a\in A, x\in\g.
$
The inclusion \(\g\hookrightarrow\e_s\) gives an inclusion \(\g\hookrightarrow\widetilde\g_s\), and the quotient is
\[
 \widetilde\g_s/\g\simeq A\otimes_BT_{B/\kk}.
\]
Define the anchor by
\[
 \widetilde\rho_s(a\otimes u)=a\sigma_s(u).
\]
This is well-defined modulo \(\cR\), since for \(x\in\g\) one has \(\sigma_s(x)=\rho(x)\), and \(\rho\) is \(A\)-linear.
We define the bracket as follows. For homogeneous sections put
\[
 [a\otimes u,b\otimes v]
 =ab\otimes[u,v]
 +a\sigma_s(u)(b)\otimes v
 -(-1)^{|u||v|}b\sigma_s(v)(a)\otimes u.
\]
Since \(A\) is concentrated in degree zero in the present affine smooth setting, no further coefficient signs occur. We verify that this formula descends to the quotient by \(\cR\). Replacing \(a\otimes x\) by \(1\otimes ax\) changes the right-hand side by the first-order Leibniz rule for the derivation of \(\g\) encoded by \(u\in\e_s\). Hence the bracket is compatible with the relations. The Jacobi identity follows by expanding the Jacobiator: the terms involving brackets in \(\e_s\) cancel by the Jacobi identity of \(\e_s\), the terms involving two symbols cancel because
\[
 \sigma_s([u,v])=[\sigma_s(u),\sigma_s(v)],
\]
and the mixed terms cancel because each \(u\in\e_s\) acts by a first-order Lie derivation of \(\g\). Therefore \((\widetilde\g_s,\widetilde\rho_s)\) is an absolute dg-Lie algebroid. Its relative inverse image is \(\g\), and its transverse quotient is \(A\otimes_BT_{B/\kk}\); hence it is a strict transversal unfolding.

\smallskip
\noindent\emph{Step 4: inverse constructions.}
Starting from an unfolding \(\widetilde\g\), define
\[
 \Psi\from\widetilde\g^{\bas}\longrightarrow
 T_{B/\kk}\times_{\uuder(\g/B)}D^1_{\bas}(\g/B),
 \qquad
 (\widetilde\xi,\xi)\longmapsto
 (\xi,\Theta_{(\widetilde\xi,\xi)}).
\]
The map is injective: if two basic lifts have the same image under \(\Psi\), their difference is an element \(x\in\g\) with \(\iota(x)=0\), and effectivity gives \(x=0\). It is surjective: given \((\xi,D)\) in the pullback, choose any basic lift \(\widehat\xi=(\widetilde\xi,\xi)\) of \(\xi\). Since \([D]=s_{\widetilde\g}(\xi)\), there is a unique \(x\in\g\) with
\[
 D-\Theta_{\widehat\xi}=\iota(x),
\]
and \(\widehat\xi+x=(\widetilde\xi+x,\xi)\) maps to \((\xi,D)\). Thus \(\Psi\) is an isomorphism. Extending this isomorphism by the \(A\)-linearisation construction recovers \(\widetilde\g\), by Lemma~\ref{lem:automatic-basic-reconstruction}. Conversely, if one starts from a splitting \(s\), the basic part of the reconstructed unfolding is \(\e_s\), and its adjoint-symbol splitting is \(s\). The two constructions are therefore inverse on isomorphism classes. The construction is functorial in morphisms of triples \((B,A,\g)\) which preserve the chosen smooth base map and the relative dg-Lie algebroid, because every step is defined by pullback, commutator, and the universal quotient by inner derivations.
\end{proof}

\subsection{Maurer--Cartan form}

Let
\[
 s_0\from T_{B/\kk}\longrightarrow\uuder(\g/B)
\]
be a flat reference splitting and put \(\kfr=\ker(a_{\uuder})\). The adjoint action of
\(s_0(T_{B/\kk})\) makes \(\kfr\) a dg-module over the Lie algebroid
\(T_{B/\kk}\). The deformation complex is the total Chevalley--Eilenberg dg-Lie algebra
\[
 \mathfrak C_{s_0}
 :=\Tot\!\left(
 C^\bullet_{\mathrm{CE}}(T_{B/\kk};\kfr)
 \right)
 \simeq
 \Tot\!\left(
 \Omega^\bullet_{B/\kk}\otimes_B\kfr
 \right).
\]
If \(\alpha\in\Omega^p_{B/\kk}\otimes_B\kfr^q\), then
\(|\alpha|_{\mathrm{tot}}=p+q\), and we use the total differential
\[
 D_{s_0}\alpha
 =d_{\mathrm{CE},\nabla^{s_0}}\alpha
   +(-1)^p d_{\kfr}\alpha.
\]
The bracket is the wedge extension of the bracket of \(\kfr\), with the standard Koszul sign. Since \(s_0\) is a morphism of dg-Lie algebroids, one has
\(D_{s_0}^2=0\).

A strict splitting differs from \(s_0\) by an element of
\(\Omega^1_{B/\kk}\otimes_B\kfr^0\). More generally, the formal derived neighbourhood of
\(s_0\) is governed by all elements of total degree one in \(\mathfrak C_{s_0}\). For such an element the flatness equation is the Maurer--Cartan equation
\[
 D_{s_0}\alpha+\frac12[\alpha,\alpha]=0.
\]
In particular, for a strict correction
\(\alpha\in\Omega^1_{B/\kk}\otimes_B\kfr^0\), the equation contains both
\[
 d_{\kfr}\alpha=0
 \qquad\text{and}\qquad
 d_{\mathrm{CE},\nabla^{s_0}}\alpha
 +\frac12[\alpha,\alpha]=0.
\]
This is the form of the Maurer--Cartan equation used throughout the deformation and obstruction theory below. Replacing the strict controller by an equivalent
\(L_\infty\)-model gives the usual complete equation
\[
 \ell_1(\alpha)+\frac12\ell_2(\alpha,\alpha)
 +\frac1{3!}\ell_3(\alpha,\alpha,\alpha)+\cdots=0
\]
in the corresponding total convolution \(L_\infty\)-algebra \cite{Hinich01,Pridham10,Nuiten}.


\section{Homotopical meaning of flat splittings}\label{sec:homotopical-controllers}

The strict affine theorem is deliberately concrete. The non-effective and intrinsic assertions belong instead to the homotopy theory of dg-Lie algebroids. We use Nuiten's Quillen-equivalent semi-model structures on dg-Lie algebroids and \(L_\infty\)-algebroids over a characteristic-zero cdga \cite{Nuiten}; their localisation will be denoted by
\[
 \mathrm{LieAlgd}^{\infty}_{B}.
\]
All mapping spaces below are derived mapping spaces in this localisation. A crossed dg-Lie algebroid is regarded, through the standard crossed-module construction, as a two-term \(L_\infty\)-algebroid. This convention distinguishes the invariant mapping space from any chosen strict presentation.

\begin{definition}\label{def:derived-flat-space}
Let \(\mathbb U\to T_{B/\kk}\) be a dg- or \(L_\infty\)-algebroid over the base tangent algebroid. Its derived space of flat splittings is
\[
 \operatorname{Flat}^{\mathrm{der}}(T_{B/\kk},\mathbb U)
 :=
 \Map_{(\mathrm{LieAlgd}^{\infty}_{B})_{/T_{B/\kk}}}
 \bigl(T_{B/\kk},\mathbb U\bigr).
\]
Thus a point is a morphism over \(T_{B/\kk}\), and therefore already includes compatibility with the internal differential, the anchor and all brackets. Higher simplices record homotopies between such morphisms.
\end{definition}

\begin{proposition}\label{prop:strict-flat-pi0}
Let \(\mathfrak u\to T_{B/\kk}\) be a strict effective controller, represented by a cofibrant-fibrant dg-Lie algebroid, and assume that the mapping space from \(T_{B/\kk}\) is connective in the range relevant to strict degree-zero sections. Then
\[
 \pi_0\operatorname{Flat}^{\mathrm{der}}(T_{B/\kk},\mathfrak u)
\]
is the ordinary set of flat dg-Lie-algebroid splittings. Under these hypotheses Theorem~\ref{thm:strict-affine} computes the connected components of the derived splitting space.
\end{proposition}

\begin{proof}
A strict point of the derived mapping space is a section of the anchor which is a morphism of complexes and preserves the Lie bracket. These are exactly the flat splittings of Definition~5.1. Effectivity removes the degree-zero stabilisers coming from the inner-action kernel, and the stated connectivity assumption ensures that no additional positive homotopy contributes to \(\pi_0\).
\end{proof}

\section{Non-effective unfoldings and the crossed-module controller}\label{sec:higher}

When the inner action has a kernel, the quotient
\(D^1_{\bas}(\g/B)/\iota(\g)\) loses genuine isotropy. The appropriate replacement is the crossed module
\[
 \Uder(\g/B)
 :=\bigl[\g\xto{\iota}D^1_{\bas}(\g/B)\bigr],
\]
where \(D^1_{\bas}(\g/B)\) acts on \(\g\) through its derivation component. The crossed-module identities are
\[
 \iota(D\cdot x)=[D,\iota(x)],
 \qquad
 \iota(x)\cdot y=[x,y].
\]
The kernel
\[
 \z(\g)=\ker(\iota)
 =\{x\in\g:\rho(x)=0,\ \operatorname{ad}_x=0\}
\]
is the pointwise central isotropy of this crossed module. Crossed modules and their associated two-groupoids provide strict presentations of two-term homotopical quotients \cite{BrownSpencer76,Noohi07}; the dg enhancement is interpreted in \(\mathrm{LieAlgd}^{\infty}_{B}\) by the preceding convention.

\begin{definition}\label{def:strict-lift-groupoid}
The strict flat-lift groupoid
\(\operatorname{Flat}_{\mathrm{str}}(T_{B/\kk},\Uder(\g/B))\)
has as its objects the basic extensions
\[
 0\longrightarrow\g\longrightarrow\e\longrightarrow T_{B/\kk}\longrightarrow0
\]
together with a morphism
\(\Theta\from\e\to D^1_{\bas}(\g/B)\) restricting to \(\iota\), inducing the identity on the base, and preserving the differential and bracket. Morphisms are isomorphisms of such extensions which are the identity on \(\g\) and on \(T_{B/\kk}\) and commute with \(\Theta\).
\end{definition}

\begin{proposition}\label{prop:strict-presentation-crossed}
For a cofibrant strictly perfect crossed presentation, the nerve of
\(\operatorname{Flat}_{\mathrm{str}}(T_{B/\kk},\Uder(\g/B))\)
identifies with the one-truncation of the derived splitting space:
\[
 N\operatorname{Flat}_{\mathrm{str}}(T_{B/\kk},\Uder(\g/B))
 \simeq
 \tau_{\leq1}
 \operatorname{Flat}^{\mathrm{der}}(T_{B/\kk},\Uder(\g/B)).
\]
The full derived mapping space retains the higher homotopies arising from the internal dg structure.
\end{proposition}

\begin{proof}
A crossed module presents a strict two-groupoid. Morphisms from the ordinary base algebroid to this presentation correspond to the extensions displayed in Definition~\ref{def:strict-lift-groupoid}, and natural transformations correspond precisely to their isomorphisms. This is the customary description of the mapping groupoid for a crossed module \cite{BrownSpencer76,Noohi07}. Passing to the dg setting and replacing source and target by admissible cofibrant-fibrant models computes the derived mapping space in Nuiten's localisation \cite{Nuiten}. The nerve of the strict mapping groupoid therefore computes its one-truncation. The higher homotopies are retained by the untruncated mapping space on the right.
\end{proof}

\begin{proposition}\label{prop:central-stabiliser-complex}
Let \(s\) be a strict flat lift and let \(\z(\g)_s\) denote the central isotropy with the action induced by \(s\). An automorphism of the corresponding strict extension is given by a degree-zero \(B\)-linear map
\[
 c\from T_{B/\kk}\longrightarrow\z(\g)
\]
which is compatible with the internal differential and satisfies
\[
 c([\xi,\eta])
 =\nabla^s_\xi c(\eta)-\nabla^s_\eta c(\xi).
\]
In particular, when the relevant complexes are concentrated in degree zero,
\[
 \operatorname{Aut}(s)
 \cong
 Z^1_{\mathrm{CE}}
 \bigl(T_{B/\kk};\z(\g)_s\bigr).
\]
In the general dg case the based loop space of
\(\operatorname{Flat}^{\mathrm{der}}\) is governed by the total complex
\[
 \Tot C^\bullet_{\mathrm{CE}}
 \bigl(T_{B/\kk};\z(\g)_s\bigr),
\]
with the customary shift determined by the two-term convention.
\end{proposition}

\begin{proof}
Choose a splitting notation for the extension and write its elements as
\(x+s(\xi)\). An automorphism which is the identity on the kernel and quotient must have the form
\[
 \Phi_c(x+s(\xi))=x+s(\xi)+c(\xi).
\]
Compatibility with the differential gives the chain condition on \(c\). Preservation of the bracket gives
\[
 c([\xi,\eta])
 =\nabla^s_\xi c(\eta)-\nabla^s_\eta c(\xi)+[c(\xi),c(\eta)].
\]
Since \(c\) takes values in the central isotropy, the last term vanishes. This is the Chevalley--Eilenberg cocycle equation. The final assertion is the derived form of the same calculation.
\end{proof}

\begin{theorem}\label{thm:higher}
Let \(B\to A\) be smooth and let
\(\g\to T_{A/B}\) be a cofibrant strictly perfect relative dg-Lie algebroid. The adjoint-symbol construction and scalar-extension reconstruction induce an equivalence of groupoids
\[
 \operatorname{Unf}^{\tr}_{\mathrm{str}}(\g/B)
 \simeq
 \operatorname{Flat}_{\mathrm{str}}
 \bigl(T_{B/\kk},\Uder(\g/B)\bigr).
\]
Consequently,
\[
 N\operatorname{Unf}^{\tr}_{\mathrm{str}}(\g/B)
 \simeq
 \tau_{\leq1}
 \operatorname{Flat}^{\mathrm{der}}
 \bigl(T_{B/\kk},\Uder(\g/B)\bigr).
\]
We define the full derived space of transversal unfoldings by the untruncated mapping space on the right. If \(\iota\) is injective, the central isotropy vanishes and the statement recovers Theorem~\ref{thm:strict-affine} on connected components.
\end{theorem}

\begin{proof}
Let \(\widetilde\g\) be a strict unfolding and let
\(\widetilde\g^{\bas}\) be its basic part. A basic lift
\(\widehat\xi=(\widetilde\xi,\xi)\) determines
\[
 \Theta_{\widehat\xi}
 =\bigl(\widetilde\rho(\widetilde\xi),
        \operatorname{ad}_{\widetilde\xi}|_\g,\xi\bigr)
 \in D^1_{\bas}(\g/B).
\]
The Lie--Rinehart identity makes this a first-order Lie derivation, and changing the lift by \(x\in\g\) changes it by \(\iota(x)\). The resulting map
\(\widetilde\g^{\bas}\to D^1_{\bas}(\g/B)\) restricts to \(\iota\), induces the identity on the base and preserves the differential and bracket. Hence it is an object of the strict flat-lift groupoid.

Conversely, an object
\((\e,\Theta)\) of that groupoid satisfies the hypotheses of Lemma~\ref{lem:reconstruction-basic-extension}. The scalar-extension construction
\[
 \widetilde\g_{\e}:=(A\otimes_B\e)/\cR
\]
therefore produces a strict transversal unfolding. The two constructions are inverse because Lemma~\ref{lem:automatic-basic-reconstruction} recovers an unfolding from its basic part, while the adjoint-symbol morphism of a reconstructed extension is the original \(\Theta\).

The same formulas identify morphisms. In particular, their stabilisers are the cocycles described in Proposition~\ref{prop:central-stabiliser-complex}, rather than isolated elements of \(\ker\iota\). The comparison with the one-truncation of the derived mapping space is Proposition~\ref{prop:strict-presentation-crossed}.
\end{proof}


\section{The intrinsic graded-mixed controller}\label{sec:intrinsic}

We now formulate the affine controller intrinsically in terms of the graded-mixed de Rham algebra of a relative derived foliation, following the intrinsic definition of derived foliations in \cite[Section~2.1]{TV}. This formulation requires no chosen Chevalley--Eilenberg presentation.

Let \(\pi\from X\to S\) be a smooth morphism of derived Artin stacks over \(\kk\), and let \(\cF/S\) be a perfect relative derived foliation in the sense of To\"en--Vezzosi \cite[Section~2.1]{TV}. Thus \(\DR(\cF/S)\) is a graded mixed cdga under \(\DR(X/S)\), locally quasi-free on its weight-one piece, and the tangent complex \(\bbT_{\cF/S}\) is dual to the weight-one cotangent complex.

\begin{definition}\label{def:basic-gm-der}
Let
\[
 R:=\DR(\cF/S).
\]
The dg-Lie algebroid of weight-zero graded-mixed derivations of \(R\) is denoted
\(\Der^{\gm}(R)\), and its symbol is restriction to the weight-zero algebra:
\[
 \sigma\from\Der^{\gm}(R)\longrightarrow\bbT_X.
\]
Restriction of a symbol to functions pulled back from the base gives a morphism of complexes
\[
 r_S\circ\sigma\from
 \Der^{\gm}(R)
 \longrightarrow
 \mathbb R\Der_{\kk}
 \bigl(\pi^{-1}\cO_S,\cO_X\bigr).
\]
The natural action of a base vector field on pulled-back functions gives
\[
 \pi^{-1}\bbT_S
 \longrightarrow
 \mathbb R\Der_{\kk}
 \bigl(\pi^{-1}\cO_S,\cO_X\bigr).
\]
We define the complex of basic graded-mixed derivations by the homotopy pullback
\[
 \Der^{\gm}_{\bas}(R)
 :=
 \Der^{\gm}(R)
 \times^h_{\mathbb R\Der_{\kk}(\pi^{-1}\cO_S,\cO_X)}
 \pi^{-1}\bbT_S.
\]
Thus a strict section is a projectable pair \((\Theta,\xi)\) satisfying
\[
 \sigma(\Theta)(\pi^{-1}f)=\pi^{-1}(\xi f)
 \qquad
 \text{for every local }f\in\cO_S.
\]
The comparison complex
\(\mathbb R\Der_{\kk}(\pi^{-1}\cO_S,\cO_X)\) is used solely to record symbols. The Lie bracket on basic derivations is defined on projectable pairs by
\[
 [ (\Theta,\xi),(\Theta',\eta) ]
 :=([\Theta,\Theta'],[\xi,\eta]).
\]
This is well-defined because the commutator of projectable derivations is projectable. In particular, no Lie bracket on the full \(\cO_X\)-module \(\pi^*\bbT_S\) is required, and no splitting of the tangent triangle is chosen.
\end{definition}

\begin{construction}\label{cons:inner-action}
After replacing \(\DR(\cF/S)\) by a cofibrant quasi-free graded mixed model, the tangent dg-Lie algebroid \(\bbT_{\cF/S}\) acts by contraction and Lie derivative; the resulting action is independent of the chosen model in the homotopy category. Consequently there is an induced morphism of dg-Lie algebroids
\[
 \iota_{\cF/S}\from \bbT_{\cF/S}\longrightarrow \Der^{\gm}_{\bas}(\DR(\cF/S)).
\]
Its symbol is the relative anchor \(\bbT_{\cF/S}\to \bbT_X\). Since this anchor factors through the fibre of \(q_X\from\bbT_X\to\pi^*\bbT_S\), its image in \(\pi^*\bbT_S\) is zero. Thus the derivation is basic, with base component equal to the zero vector field on \(S\).
\end{construction}

\begin{definition}\label{def:intrinsic-controller}
The intrinsic derived transverse controller of the relative foliation \(\cF/S\) is the homotopy quotient
\[
 \Uder(\cF/S):=
 \bigl[\bbT_{\cF/S}\xto{\iota_{\cF/S}}\Der^{\gm}_{\bas}(\DR(\cF/S))\bigr].
\]
Its invariant effective shadow is the homotopy cofibre
\[
 \uuder^{h}(\cF/S):=
 \hocofib\bigl(
 \bbT_{\cF/S}\to
 \Der^{\gm}_{\bas}(\DR(\cF/S))
 \bigr).
\]
If the inner map is represented by a cofibration which is a degreewise monomorphism, the ordinary quotient is a strict model of this homotopy cofibre.
\end{definition}

\begin{proposition}\label{prop:intrinsic-comparison}
Let \(B\to A\) be a smooth morphism between smooth commutative \(\kk\)-algebras, and let \(\cF_{A/B}\) be presented by a cofibrant strictly perfect relative dg-Lie algebroid
\[
 \rho\from\g\longrightarrow T_{A/B}
\]
through an equivalence of graded mixed cdga's
\[
 \DR(\cF_{A/B})\simeq\CE^*(\g).
\]
Then the dg-Lie algebra of basic weight-zero graded-mixed derivations of \(\CE^*(\g)\) identifies with \(D^1_{\bas}(\g/B)\). Under this identification the intrinsic inner action of \(\bbT_{\cF_{A/B}/B}\) is the map \(\iota\from\g\to D^1_{\bas}(\g/B)\). Hence the intrinsic controller is represented by the crossed module \([\g\to D^1_{\bas}(\g/B)]\), and by the ordinary quotient \(\uuder(\g/B)\) in the effective case.
\end{proposition}

\begin{proof}
We include the calculation because it identifies the graded-mixed formulation with the Lie--Rinehart formulation.

Set
 $
 R:=\CE^*(\g)=\Sym_A(\g^\vee[1]).
$
Throughout this calculation we adopt the convention that \(\CE^*(\g)=\Sym_A(\g^\vee[1])\), the mixed differential has bidegree \((-1,+1)\), and derivations are cohomological derivations of degree \(n\). The algebra \(R\) is free as a graded commutative algebra over \(A\), generated in weight one by \(\g^\vee[1]\). Hence a homogeneous weight-zero graded derivation \(\Theta\) of degree \(n\) of the underlying graded cdga is determined by the following two components:
\[
 \Theta_0:=\Theta|_A\from A\longrightarrow A[n]
\]
and
\[
 \Theta_1:=\Theta|_{\g^\vee[1]}\from
 \g^\vee[1]\longrightarrow \g^\vee[1][n],
\]
subject only to the graded Leibniz rule. The first component is a derivation \(\theta\in T_{A/\kk}^n\). Since \(\g\) is perfect, the second component may be dualised and desuspended; it is equivalently a \(\kk\)-linear operator
\[
 \delta\from\g\longrightarrow\g[n]
\]
satisfying the first-order Leibniz rule
\[
 \delta(ax)=\theta(a)x+(-1)^{n|a|}a\delta(x)
\]
for homogeneous \(a\in A\) and \(x\in\g\). Thus, before imposing compatibility with the mixed differential, a homogeneous weight-zero graded derivation of \(R\) is equivalent to a pair \((\theta,\delta)\), where \(\delta\) is a first-order operator on \(\g\) with symbol \(\theta\).

The graded mixed algebra \(R\) has two differentials. The internal cohomological differential is induced by the differential of the dg-Lie algebroid \(\g\). The mixed Chevalley--Eilenberg differential, denoted \(\epsilon_{CE}\), is the sum of the component dual to the anchor and the component dual to the bracket. A graded mixed derivation is a graded derivation commuting with \(\epsilon_{CE}\); the cohomological differential on the derivation complex is the commutator with the internal differential. Thus the structural condition is
\[
 [\Theta,\epsilon_{CE}]=0,
\]
not the assertion that \(\Theta\) is a closed element of the total derivation complex. We record the two substantive components of this equation.

First, evaluate the commutator on functions. For \(a\in A\), the anchor component of \(\epsilon_{CE}\) sends \(a\) to the functional \(x\mapsto \rho(x)(a)\). The identity \([\Theta,\epsilon_{CE}](a)=0\) says that, for every \(x\in\g\),
\[
 \theta(\rho(x)(a))-
\rho(\delta x)(a)-(-1)^{|\Theta||x|}\rho(x)(\theta(a))=0.
\]
Equivalently,
\[
 \rho(\delta x)=[\theta,\rho(x)] .
\]
Thus the symbol \(\theta\) and the operator \(\delta\) are compatible with the anchor.

Second, evaluate the commutator on a linear generator \(\alpha\in\g^\vee\). The bracket part of \(\epsilon_{CE}\alpha\) is dual to the Lie bracket on \(\g\). The identity \([\Theta,\epsilon_{CE}](\alpha)=0\), evaluated on homogeneous \(x,y\in\g\), becomes
\[
 \delta([x,y])=[\delta x,y]+(-1)^{|\delta||x|}[x,\delta y].
\]
Thus \(\delta\) is a derivation of the dg-Lie bracket. The internal differential of \(\g\) is not an additional algebraic constraint on the underlying graded derivation; it is the differential of the derivation complex, sending \((\theta,\delta)\) to the graded commutator with the internal differential. Under the correspondence above, this is the differential of \(D^1_{A/\kk}(\g)\).

Consequently graded-mixed derivations of \(\CE^*(\g)\) are pairs \((\theta,\delta)\) satisfying the structural conditions in the definition of \(D^1_{A/\kk}(\g)\), with the same cohomological differential. The basic condition is also read from the transverse symbol: \(\Theta\) is basic when the image of \(\theta\) under
\[
 T_{A/\kk}\longrightarrow A\otimes_B T_{B/\kk}
\]
lies in \(1\otimes T_{B/\kk}\). Together with the chosen base vector field \(\xi\), this identifies a basic graded-mixed derivation with the triple \((\theta,\delta,\xi)\) defining \(D^1_{\bas}(\g/B)\). Hence
\[
 \Der^{\gm}_{\bas}(\CE^*(\g))\cong D^1_{\bas}(\g/B)
\]
as dg-Lie algebras. The bracket is the commutator on both sides: on the CE side it is the commutator of derivations of the graded mixed cdga, and on the Lie--Rinehart side it is the commutator of first-order Lie derivations.

It remains to compare the inner actions. A tangent element \(x\in\g\) acts on \(R\) by the usual Cartan Lie derivative \(L_x=[\epsilon_{CE},\iota_x]\). Under the preceding description this derivation has symbol \(\rho(x)\) and acts on \(\g\) by \(\operatorname{ad}_x\). Hence it corresponds exactly to
\[
 \iota(x)=(\rho(x),\operatorname{ad}_x,0)\in D^1_{\bas}(\g/B).
\]
Taking the homotopy quotient by these inner derivations gives the crossed module \([\g\to D^1_{\bas}(\g/B)]\). If \(\iota\) is injective and is represented by a cofibration which is a degreewise monomorphism, the ordinary quotient \(\uuder(\g/B)\) is a strict model of the corresponding homotopy cofibre.
\end{proof}

We now pass from the affine classification to the global statement, under explicit hypotheses of CE-presentability and descent. The argument uses the descent theory of derived foliations on derived stacks developed in \cite[Section~2.1]{TV}, Illusie's treatment of exact cotangent triangles under base change \cite{IllusieI}, and homotopical descent as in \cite{HAGII}.

\begin{definition}\label{def:local-groupoids}
Let \(U_\bullet\to X\) be a smooth affine hypercover carrying cofibrant strictly perfect CE presentations. For each \(n\), let
\[
 \mathcal U_n^{\mathrm{der}}
\]
be the derived local unfolding space defined by Theorem~\ref{thm:higher}, and let
\[
 \mathcal S_n^{\mathrm{der}}
 :=\operatorname{Flat}^{\mathrm{der}}
 \bigl(T_S,\Uder(\cF|_{U_n}/S)\bigr).
\]
Their one-truncations are the nerves of the strict local unfolding and flat-lift groupoids. Pullback along every simplicial morphism of the hypercover gives maps of spaces, so
\((\mathcal U_n^{\mathrm{der}})_n\) and
\((\mathcal S_n^{\mathrm{der}})_n\) are cosimplicial derived spaces.
\end{definition}

\begin{hypothesis}\label{hyp:global-descent}
Let \(\pi\from X\to S\) be a smooth morphism of derived Artin stacks over \(\kk\), and let \(\cF/S\) be a perfect relative derived foliation. We assume:
\begin{enumerate}[label=(\roman*)]
\item there is a smooth affine hypercover \(U_\bullet\to X\) such that each \(\cF|_{U_n}/S\) admits a cofibrant strictly perfect CE presentation;
\item the affine classification is functorial with respect to all face and degeneracy morphisms of the hypercover;
\item local unfoldings and local derived splitting spaces satisfy smooth descent;
\item the local crossed-module controllers descend to a stack of higher Lie algebroids \(\Uder(\cF/S)\) over \(X\);
\item the derived pushforward of this controller exists in the chosen \(\infty\)-category of Lie-algebroid stacks over \(S\), satisfies smooth hyperdescent, and carries a descended anchor
\[
 a\from \mathrm R\pi_*\Uder(\cF/S)\longrightarrow \bbT_S.
\]
\end{enumerate}
In the effective case we write \(\uuder(\cF/S)\) for the one-categorical truncation of \(\Uder(\cF/S)\).
\end{hypothesis}

\begin{proposition} \label{prop:checkable-global-descent}
Hypothesis~\ref{hyp:global-descent} holds whenever the following concrete conditions are satisfied for a smooth affine hypercover \(U_\bullet\to X\):
\begin{enumerate}[label=(\alph*)]
\item each \(\cF|_{U_n}/S\) is represented by a cofibrant strictly perfect CE dg-Lie algebroid \(\g_n\), and the induced local CE controllers form a cosimplicial crossed controller satisfying smooth descent;
\item the vertical complexes of the local controllers are perfect and \(\mathrm R\pi_*\)-perfect, so that their derived pushforwards exist in the chosen \(\infty\)-category of controller stacks over \(S\);
\item the bracket, anchor, inner action and crossed-module action descend under \(\mathrm R\pi_*\), and the descended anchor
\[
 a\from \mathrm R\pi_*\Uder(\cF/S)\longrightarrow \bbT_S
\]
is compatible with the cosimplicial descent data.
\end{enumerate}
Then the local derived classification glues to the global classification of Theorem~\ref{thm:intrinsic}. The criterion does not purport to construct a general theory of pushforwards for Lie-algebroid stacks; it isolates the precise descent hypotheses required here.
\end{proposition}

\begin{proof}
The proof is an application of descent in the chosen \(\infty\)-category of controller stacks. By (a), the local CE controllers form a cosimplicial crossed controller and the local Lie--Rinehart identities descend as equalities of morphisms of perfect complexes. By Proposition~\ref{prop:intrinsic-comparison}, these local CE controllers present the restrictions of the intrinsic object
\[
 \Uder(\cF/S)=
 \bigl[\bbT_{\cF/S}\to\Der^{\gm}_{\bas}(\DR(\cF/S))\bigr].
\]
Condition (b) supplies the required derived pushforward, and condition (c) ensures that the pushed-forward bracket, anchor and crossed-module structure remain compatible with descent.

Flat splittings are regarded not merely as a set of sections, but as the derived mapping space
\[
 \operatorname{Flat}^{\mathrm{der}}_S(\bbT_S,\Uder(\cF/S))
 :=
 \Map_{\mathrm{LieAlgdStack}_{/\bbT_S}}
 \bigl(\bbT_S,\mathrm R\pi_*\Uder(\cF/S)\bigr),
\]
or the analogous crossed mapping space in the non-effective case. Mapping spaces in an \(\infty\)-category satisfy descent. Therefore the local affine equivalences glue by taking homotopy limits over the hypercover.
\end{proof}

\begin{theorem}\label{thm:intrinsic}
Under Hypothesis~\ref{hyp:global-descent}, transversal derived unfoldings of \(\cF/S\) are equivalent to the derived space of flat splittings of the descended controller:
\[
 \operatorname{Unf}^{\tr}_{\mathrm{der}}(\cF/S)
 \simeq
 \operatorname{Flat}^{\mathrm{der}}_S(\bbT_S,\Uder(\cF/S))
 :=
 \Map_{\mathrm{LieAlgdStack}_{/\bbT_S}}
 \bigl(\bbT_S,\mathrm R\pi_*\Uder(\cF/S)\bigr),
\]
where the target is computed in the localised \(L_\infty\)-algebroid stack. Its one-truncation classifies strict descended unfoldings. In the effective strict case this equivalence induces a bijection on connected components with ordinary flat splittings of \(\mathrm R\pi_*\uuder(\cF/S)\to\bbT_S\).
\end{theorem}

\begin{proof}
For every \(n\), Theorem~\ref{thm:higher} and its effective specialisation provide an equivalence of derived spaces
\[
 \Phi_n\from \mathcal U_n^{\mathrm{der}}
 \longrightarrow \mathcal S_n^{\mathrm{der}}.
\]
By assumption~(ii), these equivalences are compatible with all coface and codegeneracy morphisms. Thus
\((\Phi_n)_n\) is a levelwise equivalence of cosimplicial spaces. Taking the homotopy limit over \(\Delta\) gives
\[
 \operatorname*{holim}_{[n]\in\Delta}
 \mathcal U_n^{\mathrm{der}}
 \xrightarrow{\ \sim\ }
 \operatorname*{holim}_{[n]\in\Delta}
 \mathcal S_n^{\mathrm{der}}.
\]
Assumption~(iii) identifies the left-hand side with the global derived unfolding space. The right-hand side is the mapping space into the descended pushed-forward controller, because mapping spaces in an \(\infty\)-category satisfy hyperdescent. This proves
\[
 \operatorname{Unf}^{\tr}_{\mathrm{der}}(\cF/S)
 \simeq
 \operatorname{Flat}^{\mathrm{der}}_S
 \bigl(\bbT_S,\Uder(\cF/S)\bigr).
\]
Taking one-truncations yields the strict descended groupoid; in the effective strict case, passing further to connected components recovers ordinary flat splittings.
\end{proof}

\begin{remark}
Theorem \ref{thm:intrinsic} is stated as a descent theorem. The affine classification provides the local theorem; the global statement is obtained by gluing, and no stronger claim is made. This formulation retains the intrinsically homotopical nature of the quotient by inner symmetries.
\end{remark}


\section{Homotopy invariance of the controller}\label{sec:homotopy-invariance}

The naively defined derivation complex of a non-cofibrant algebra is not homotopy invariant. The invariant replacement is the tangent \(L_\infty\)-algebroid of the formal moduli problem of projectable graded-mixed automorphisms. Although a bimodule of derivations may compare the underlying complexes, it does not by itself carry the requisite strict bracket.

\begin{definition}
Let \(R\) be a cofibrant quasi-free graded-mixed cdga under \(\DR(A/B)\). Denote by
\(\mathcal{A}ut^{\mathrm{gm}}_{\bas}(R)\) the formal moduli problem which assigns to a nilpotent augmented cdga \(C\) the space of graded-mixed automorphisms of
\(R\otimes C\) reducing to the identity modulo the augmentation and whose symbols are projectable over the base. Its tangent \(L_\infty\)-algebroid is denoted
\[
 \mathbb R\Der^{\mathrm{gm}}_{\bas}(R).
\]
For a cofibrant quasi-free model its underlying tangent complex is the projectable derivation complex of Definition~\ref{def:basic-gm-der}, and its binary bracket is the commutator of derivations.
\end{definition}

The characteristic-zero equivalence between formal moduli problems and dg- or \(L_\infty\)-algebroids is used here in its relative form; see \cite{Hinich01,Pridham10,Nuiten,NuitenKoszul}. The semi-model structures of \cite{Nuiten} provide admissible cofibrant replacements, while the formal-moduli interpretation determines the invariant higher brackets.

\begin{proposition}\label{prop:homotopy-invariance-der}
Let \(R\to R'\) be an equivalence of perfect cofibrant quasi-free graded-mixed cdga's under \(\DR(A/B)\). Conjugation along this equivalence induces an equivalence of formal moduli problems
\[
 \mathcal{A}ut^{\mathrm{gm}}_{\bas}(R)
 \simeq
 \mathcal{A}ut^{\mathrm{gm}}_{\bas}(R'),
\]
and hence an equivalence of tangent \(L_\infty\)-algebroids
\[
 \mathbb R\Der^{\mathrm{gm}}_{\bas}(R)
 \simeq
 \mathbb R\Der^{\mathrm{gm}}_{\bas}(R').
\]
The equivalence is compatible with the projectable symbol and with the tangent-inner Cartan action.
\end{proposition}

\begin{proof}
An equivalence \(R\to R'\) induces, after derived base change to every nilpotent augmented test algebra, an equivalence between the spaces of automorphisms reducing to the identity. The projectability condition is expressed by the homotopy pullback of symbol complexes in Definition~\ref{def:basic-gm-der} and is therefore preserved. This gives an equivalence of the two formal moduli problems. Passing to tangent \(L_\infty\)-algebroids yields the displayed equivalence by relative formal deformation theory \cite{Hinich01,Pridham10,NuitenKoszul}.

The tangent-inner map is obtained from the Cartan formula
\(L_x=[\epsilon,\iota_x]\). Contraction and Lie derivative are natural under equivalences of cofibrant graded-mixed algebras, so the equivalence of tangent algebroids intertwines the inner actions and their anchors.
\end{proof}

\begin{theorem}\label{thm:presentation-independence}
Let \(\cF_{A/B}\) be a perfect relative derived foliation and suppose that
\[
 \CE^*(\g)\simeq\DR(\cF_{A/B})\simeq\CE^*(\g')
\]
are two cofibrant strictly perfect CE presentations. Then their crossed controllers are equivalent as two-term \(L_\infty\)-algebroids over \(T_{B/\kk}\):
\[
 \bigl[\g\to D^1_{\bas}(\g/B)\bigr]
 \simeq
 \bigl[\g'\to D^1_{\bas}(\g'/B)\bigr].
\]
Both represent the intrinsic controller \(\Uder(\cF_{A/B})\). In the effective case the invariant quotient is the homotopy cofibre
\[
 \hocofib\bigl(\g\to D^1_{\bas}(\g/B)\bigr),
\]
and the analogous homotopy cofibre for \(\g'\). If the inner maps are cofibrations represented by degreewise monomorphisms, the ordinary quotients \(\uuder(\g/B)\) and \(\uuder(\g'/B)\) are strict models of these equivalent homotopy cofibres.
\end{theorem}

\begin{proof}
By Proposition~\ref{prop:intrinsic-comparison}, the strict basic derivation algebras of the two CE models represent the tangent algebroids
\(\mathbb R\Der^{\mathrm{gm}}_{\bas}(\CE^*(\g))\) and
\(\mathbb R\Der^{\mathrm{gm}}_{\bas}(\CE^*(\g'))\). Proposition~\ref{prop:homotopy-invariance-der} identifies these tangent \(L_\infty\)-algebroids and intertwines the two tangent-inner maps. Taking the homotopy quotient by this natural inner action gives the equivalence of crossed controllers. Homotopy cofibres are invariant under equivalence. Under the final cofibration hypothesis, the ordinary cokernel computes the homotopy cofibre in the chosen strict semi-model presentation.
\end{proof}

\begin{corollary}\label{cor:classification-invariance}
The derived splitting space
\[
 \operatorname{Flat}^{\mathrm{der}}
 \bigl(T_{B/\kk},\Uder(\cF_{A/B})\bigr)
\]
and its one-truncation classifying strict unfoldings are independent of the chosen cofibrant CE presentation.
\end{corollary}

\begin{proof}
Equivalent targets have equivalent derived mapping spaces in an \(\infty\)-category. Apply Theorem~\ref{thm:presentation-independence} and then take the one-truncation when strict unfoldings are desired.
\end{proof}


\section{Crystals and the Gauss--Manin connection}\label{sec:crystals}

Let \(R=\DR(\cF/S)\). The coefficient systems used below are quasi-coherent crystals in the sense of To\"en--Vezzosi, rather than arbitrary graded-mixed modules.

\begin{definition}\label{def:crystals}
The \(\infty\)-category \(\Crys(\cF/S)\) is the full subcategory of graded-mixed \(R\)-modules \(M\) such that
\begin{enumerate}[label=(\roman*)]
\item \(M^{(0)}\) is quasi-coherent over \(\cO_X\); and
\item the canonical morphism of underlying graded modules
\[
 M^{(0)}\otimes_{\cO_X}^{\mathbb L}R^{\mathrm{gr}}
 \longrightarrow M^{\mathrm{gr}}
\]
is an equivalence.
\end{enumerate}
Perfect crystals are those for which \(M^{(0)}\) is perfect. The foliated de Rham complex is the Tate realisation
\[
 \DR_{\cF/S}(M):=M^{\Tate},
\]
in accordance with the graded-mixed realisation of \cite{TV,TVRH}.
\end{definition}

\begin{definition}\label{def:linearised-crystal}
Let
\[
 s\in\operatorname{Flat}^{\mathrm{der}}_S
 \bigl(\bbT_S,\Uder(\cF/S)\bigr)
\]
be a transverse splitting. A \emph{coherently \(s\)-Cartan-linearised crystal} is a crystal
\(M\in\Crys(\cF/S)\) endowed with a representation of the semidirect
\(L_\infty\)-algebroid determined by the splitting \(s\), together with a homotopy-coherent Cartan null-homotopy of the tangent-inner action. In a strict effective CE chart this consists of lifted Lie derivatives
\[
 L^M_{\Theta_\xi}\from M\longrightarrow M
\]
for representatives \(\Theta_\xi\) of \(s(\xi)\), and contractions \(\iota_v\) satisfying
\[
 [d_{\DR_{\cF/S}(M)},\iota_v]
 =L_v-\iota_{d_{\bbT}v},
\]
together with the bracket, module and descent coherences. In the non-effective case the same datum includes the homotopies associated with the central component of the crossed controller.
\end{definition}

\begin{construction}\label{cons:GM-local-action}
In a strict CE chart \(R=\CE^*(\g)\), choose a representative
\(\Theta_\xi\in D^1_{\bas}(\g/B)\) of the transverse class of a local vector field \(\xi\). The linearisation supplies a derivation
\[
 L^M_{\Theta_\xi}\from M\longrightarrow M
\]
over \(\Theta_\xi\). Its Tate realisation acts on \(\DR_{\cF/S}(M)\). The full Cartan-linearisation records coherent homotopies under changes of representative, including those measured by the crossed-module isotropy.
\end{construction}

\begin{lemma}\label{lem:cartan-homotopy}
For a coherently Cartan-linearised crystal, the tangent-inner action is null-homotopic on the foliated de Rham complex. More precisely, a closed degree-zero tangent element \(v\) satisfies
\[
 L_v=[d_{\DR_{\cF/S}(M)},\iota_v],
\]
and the corresponding statement for a non-closed or higher tangent homotopy is supplied by the full Cartan identity in Definition~\ref{def:linearised-crystal}.
\end{lemma}

\begin{proof}
This identity is part of the Cartan-linearisation structure. In a strict CE presentation it is the usual Cartan formula on the Chevalley--Eilenberg module. The prescribed higher coherences make these local null-homotopies compatible with changes of representative and with descent.
\end{proof}

\begin{theorem}\label{thm:GM}
Let \(\pi\from X\to S\) be smooth, let \(\cF/S\) satisfy the hypotheses of Theorem~\ref{thm:intrinsic}, and let
\[
 s\in\operatorname{Flat}^{\mathrm{der}}_S
 \bigl(\bbT_S,\Uder(\cF/S)\bigr).
\]
If \(M\) is a coherently \(s\)-Cartan-linearised crystal and
\(R\pi_*\DR_{\cF/S}(M)\) exists, then this derived pushforward carries a canonical flat Gauss--Manin connection, equivalently a crystal structure over the de Rham foliation of \(S\):
\[
 \nabla^{\GM}_s\from
 R\pi_*\DR_{\cF/S}(M)
 \longrightarrow
 \bbL_S\otimes_{\cO_S}
 R\pi_*\DR_{\cF/S}(M).
\]
In a strict effective CE model the connection is represented by the ordinary operators
\[
 \nabla^{\GM}_{s,\xi}
 =R\pi_*\bigl(L^M_{\Theta_\xi}\bigr).
\]
\end{theorem}

\begin{proof}
Work first on a CE hypercover. A representative \(\Theta_\xi\) of a base vector field acts on the relative foliated de Rham complex by the operator
\(L^M_{\Theta_\xi}\). Its symbol is \(\xi\), so the module derivation identity gives
\[
 L^M_{\Theta_\xi}(f\alpha)
 =\xi(f)\alpha+fL^M_{\Theta_\xi}(\alpha)
\]
for local \(f\in\cO_S\). These operators therefore define a connection after derived pushforward.

Changing a representative changes the operator by a tangent-inner action and, in the non-effective case, by the accompanying higher isotropy homotopy. Definition~\ref{def:linearised-crystal} and Lemma~\ref{lem:cartan-homotopy} identify these changes by canonical homotopies. Hence the connection depends only on the point \(s\) of the derived splitting space.

The curvature is represented locally by
\[
 [L^M_{\Theta_\xi},L^M_{\Theta_\eta}]
 -L^M_{\Theta_{[\xi,\eta]}}.
\]
The Maurer--Cartan equation for the derived splitting identifies this curvature with a tangent-inner action together with its prescribed higher null-homotopy. The coherent Cartan action therefore renders the curvature null-homotopic in the derived category. Hence the connection is flat in the sense of crystals.

All operators and homotopies are functorial under restriction along the chosen hypercover. They satisfy the descent identities because the linearisation is a representation of the descended semidirect \(L_\infty\)-algebroid. The local connections consequently glue to the asserted global Gauss--Manin connection.
\end{proof}

\begin{remark}
The Cartan hypothesis is essential. A relative crystal provides leafwise transport, while the splitting provides transverse Lie derivatives. Independence of representatives and flatness require, in addition, coherent contractions for tangent-inner actions; in the non-effective case these contractions must include the central higher homotopies of the crossed controller.
\end{remark}


\section{\'Etale base change and descent}\label{sec:base-change}

The affine controller is functorial under \'etale change of the base. The proof must use the full Atiyah sequence, or equivalently the first principal-parts complex, before taking the fibre over a prescribed symbol: operators with a fixed symbol form an affine torsor, not a module.

\begin{definition}
Let \(E\) be a perfect \(A\)-complex. Write
\(\operatorname{At}_{A/\kk}(E)\) for the derived Atiyah complex of first-order operators on \(E\), with symbol morphism
\[
 \sigma_E\from\operatorname{At}_{A/\kk}(E)\longrightarrow T_{A/\kk}.
\]
For \(\theta\in T_{A/\kk}\), the space of operators with symbol \(\theta\) is the homotopy fibre
\[
 \operatorname{Diff}^1_\theta(E,E)
 :=\operatorname{hofib}_\theta(\sigma_E).
\]
It is an affine torsor over \(\RHom_A(E,E)\), rather than an \(A\)-module.
\end{definition}

\begin{lemma}\label{lem:first-order-etale-base-change}
Let \(B\to C\) be \'etale, put \(A_C=A\otimes_BC\), and set
\(E_C=A_C\otimes_A^{\mathbb L}E\). Then first principal parts and Atiyah complexes commute with base change:
\[
 A_C\otimes_A^{\mathbb L}\operatorname{At}_{A/\kk}(E)
 \simeq
 \operatorname{At}_{A_C/\kk}(E_C).
\]
The equivalence intertwines symbol maps. Consequently, for every symbol \(\theta\) and its base change \(\theta_C\), there is an equivalence of affine derived fibres
\[
 \operatorname{Diff}^1_\theta(E,E)\times_{\Spec A}\Spec A_C
 \simeq
 \operatorname{Diff}^1_{\theta_C}(E_C,E_C).
\]
\end{lemma}

\begin{proof}
The Atiyah complex is dual to the first principal-parts extension of \(E\). Since \(A\to A_C\) is \'etale, its relative cotangent complex vanishes and the transitivity triangle identifies
\(\bbL_{A_C/\kk}\) with
\(A_C\otimes_A^{\mathbb L}\bbL_{A/\kk}\). The first principal-parts extension therefore base-changes to the corresponding extension over \(A_C\). Perfectness of \(E\) permits dualisation and gives the equivalence of Atiyah complexes. Homotopy fibres commute with this flat derived base change, which proves the assertion for prescribed symbols.
\end{proof}

\begin{proposition}\label{prop:base-change-D}
Assume \(\g\) is perfect over \(A\). For every \'etale morphism \(B\to C\), with
\(\g_C=A_C\otimes_A\g\), there is a natural equivalence of dg-Lie algebras
\[
 C\otimes_B^{\mathbb L}D^1_{\bas}(\g/B)
 \simeq
 D^1_{\bas}(\g_C/C),
\]
compatible with the symbol, the bracket and the inner map. Hence
\[
 C\otimes_B^{\mathbb L}\Uder(\g/B)
 \simeq
 \Uder(\g_C/C)
\]
as crossed controllers.
\end{proposition}

\begin{proof}
The underlying first-order operators base-change by Lemma~\ref{lem:first-order-etale-base-change}. Preservation of the dg-Lie bracket and compatibility with the anchor are homotopy-fibre conditions on morphisms between perfect complexes, hence commute with flat base change. Since
\(C\otimes_BT_{B/\kk}\simeq T_{C/\kk}\), the projectability condition on the symbol also base-changes. The commutator bracket and the inner map are functorial, giving the asserted equivalence of crossed controllers.
\end{proof}

\begin{corollary}\label{cor:base-change-unf}
Derived flat splitting spaces and strict unfolding groupoids commute with \'etale base change. In the effective case the same holds for the ordinary quotient whenever it is a strict model of the homotopy cofibre.
\end{corollary}

\begin{proof}
Apply Proposition~\ref{prop:base-change-D} and the functoriality of derived mapping spaces. Strict unfoldings commute with base change because the fibre-product basic part and the scalar-extension reconstruction are formed from operations compatible with flat descent.
\end{proof}

\begin{theorem}\label{thm:descent-controller}
Let \(\pi\from X\to S\) be smooth, with \(X\) and \(S\) smooth schemes, and let \(\cF/S\) be locally strictly CE-presented by perfect models. On the small \'etale site of \(S\), the intrinsic crossed controllers form a stack of \(L_\infty\)-algebroids, and their derived flat splitting spaces satisfy \'etale descent. The one-truncations classify the descended strict unfoldings.
\end{theorem}

\begin{proof}
Choose an \'etale affine cover carrying cofibrant perfect CE presentations. Proposition~\ref{prop:base-change-D} identifies the controllers on double overlaps, and functoriality gives the cocycle coherence on higher overlaps. Thus the local controllers descend in the localised \(\infty\)-category. Mapping stacks into a descended object satisfy descent, so the derived splitting spaces glue. The reconstruction of strict unfoldings uses fibre products, perfect tensor products and the relation module \(\cR\), all compatible with \'etale descent. The resulting strict groupoid is the one-truncation of the descended derived mapping space by Proposition~\ref{prop:strict-presentation-crossed}.
\end{proof}

\begin{remark}
For a non-\'etale change of base, a vector field on the new base need not be pulled back from the old one. The comparison then involves the relative tangent complex of the base morphism; no strict base-change assertion of the preceding form is intended.
\end{remark}


\section{The deformation complex of a fixed unfolding}\label{sec:def-complex}

Fix a flat splitting
\(s\from T_{B/\kk}\to\uuder(\g/B)\) in the effective case and let
\[
 \kfr_s:=\ker(a_{\uuder}).
\]
The adjoint action
\(\nabla^s_\xi(k)=[s(\xi),k]\) is a flat dg connection. The relevant deformation complex is the total Chevalley--Eilenberg dg-Lie algebra
\[
 \mathfrak C_s
 :=\Tot C^\bullet_{\mathrm{CE}}
 \bigl(T_{B/\kk};\kfr_s\bigr)
 \simeq
 \Tot\bigl(\Omega^\bullet_{B/\kk}\otimes_B\kfr_s\bigr).
\]
For \(\alpha\in\Omega^p_{B/\kk}\otimes_B\kfr_s^q\), put
\[
 |\alpha|_{\mathrm{tot}}=p+q,
 \qquad
 D_s\alpha
 =d_{\mathrm{CE},\nabla^s}\alpha
  +(-1)^p d_{\kfr_s}\alpha.
\]
The bracket is
\[
 [\omega\otimes x,\eta\otimes y]
 =(-1)^{|x|\deg(\eta)}
  (\omega\wedge\eta)\otimes[x,y].
\]
The flatness of \(s\) implies \(D_s^2=0\).

\begin{proposition}\label{prop:MC-def}
The formal derived deformation problem of the unfolding \(s\) is governed by \(\mathfrak C_s\). Its Maurer--Cartan elements are the total degree-one solutions of
\[
 D_s\alpha+\frac12[\alpha,\alpha]=0.
\]
A strict deformation of the splitting is represented by a solution whose leading component belongs to
\(\Omega^1_{B/\kk}\otimes_B\kfr_s^0\). The tangent and primary obstruction spaces are
\[
 T_s\operatorname{Unf}^{\tr}\simeq H^1(\mathfrak C_s),
 \qquad
 \operatorname{Obs}_s\subset H^2(\mathfrak C_s).
\]
\end{proposition}

\begin{proof}
A strict correction \(s+\alpha\) has the same anchor precisely when \(\alpha\) is vertical. Compatibility with the internal differential is the component
\(d_{\kfr_s}\alpha=0\), while preservation of the bracket is the horizontal curvature equation. Together these are exactly the total Maurer--Cartan equation. Allowing all total degree-one components gives the derived deformation problem. Linearisation gives the cocycle equation \(D_s\alpha_1=0\), while gauge equivalence through total degree-zero elements changes \(\alpha_1\) by a boundary. The standard recursive Maurer--Cartan calculation places the obstruction to continuation in total degree two \cite{Hinich01,Pridham10}.
\end{proof}

\begin{corollary}\label{cor:rigidity}
If \(H^1(\mathfrak C_s)=0\), the unfolding is infinitesimally rigid. If \(H^2(\mathfrak C_s)=0\), every nilpotent infinitesimal deformation extends formally, subject to the usual completeness hypothesis on the Maurer--Cartan filtration.
\end{corollary}

\begin{proof}
The assertion is the usual tangent-obstruction consequence of Proposition~\ref{prop:MC-def}.
\end{proof}

\begin{remark}
For the crossed controller, the formal neighbourhood of a splitting is governed by the total convolution \(L_\infty\)-algebra obtained by twisting
\(\Map(T_{B/\kk},\Uder(\g/B))\) at that splitting. Its Maurer--Cartan equation includes the higher components associated with central isotropy. The effective dg-Lie algebra \(\mathfrak C_s\) is its strict quotient when the inner action is injective.
\end{remark}

\subsection{The first-order Suwa term}

Let the base direction be generated by \(\partial_t\), and let a codimension-one foliation be generated locally by an integrable one-form \(\omega\). An unfolded generator can be written to first order as
\[
 \widetilde\omega=\omega+t\eta+h\,dt.
\]
If \(\widetilde Y=\partial_t+Y\) is a lifted parameter direction tangent to the unfolded foliation, then at \(t=0\)
\[
 h=-\omega(Y).
\]
Thus \(h\) is the local coefficient of the adjoint-symbol class of the lift.

\begin{proposition}\label{prop:suwa-class}
The coefficient \(h\) is a local representative of the degree-zero transverse-controller class associated with \(\partial_t\). If the conormal generator is changed to
\(\omega'=u\omega\), then the unfolded generator may be chosen so that
\[
 h'=u h.
\]
If the lifted vector field is changed by a vector field \(Z\) tangent to the original foliation, then
\(\omega(Z)=0\) and the coefficient \(-\omega(Y)\) is unchanged. Hence the collection of local coefficients transforms with the conormal line and represents the same controller class. The coefficient of \(dt\) in
\(\widetilde\omega\wedge d\widetilde\omega\) is the linearised Maurer--Cartan equation for this class.
\end{proposition}

\begin{proof}
The identity \(h=-\omega(Y)\) follows from
\(\widetilde\omega(\partial_t+Y)=0\) at \(t=0\). Multiplication of the unfolded generator by a unit with constant term \(u\) multiplies both \(\omega\) and \(h\) by \(u\); terms involving \(t\,dt\) vanish to first order. Replacing \(Y\) by \(Y+Z\), with \(Z\) tangent, leaves \(-\omega(Y)\) unchanged. These are precisely the changes of generator and tangent representative by which the adjoint-symbol class is defined. Finally, expanding the Frobenius equation and taking the \(dt\)-component gives the condition that the lifted first-order derivation preserve the Pfaff ideal. This is the degree-one cocycle equation in the total Chevalley--Eilenberg complex.
\end{proof}


\section{The Atiyah--Kodaira--Spencer class}\label{sec:AKS}

The anchor sequence of the controller is the derived analogue of an Atiyah sequence \cite{Atiyah57}; its connecting morphism is the first obstruction to choosing transverse directions. Curvature is a second, nonlinear obstruction which must be formed in the total Chevalley--Eilenberg complex.

\begin{definition}
Assume that the effective controller exists over \(S\), and set
\[
 \mathfrak K_{\cF/S}
 :=\fib\left(
 \pi_*\uuder(\cF/S)\longrightarrow\bbT_S
 \right).
\]
The exact triangle
\[
 \mathfrak K_{\cF/S}\longrightarrow
 \pi_*\uuder(\cF/S)\longrightarrow
 \bbT_S\xto{\kappa_{\cF/S}}
 \mathfrak K_{\cF/S}[1]
\]
defines the Atiyah--Kodaira--Spencer morphism. In the non-effective theory the same definition uses the homotopy fibre of the crossed-controller anchor.
\end{definition}

\begin{proposition}\label{prop:AKS-splitting}
The class \(\kappa_{\cF/S}\) vanishes if and only if the anchor admits a splitting in the derived category. On an acyclic cover on which the controller is represented by perfect complexes, it is represented by the \v{C}ech cocycle of differences of local derived splittings.
\end{proposition}

\begin{proof}
This is the usual splitting criterion for an exact triangle. The \v{C}ech representative is obtained by choosing local null-homotopies of the connecting morphism; their differences lie in the homotopy fibre and satisfy the cocycle identity.
\end{proof}

Let \(\sigma\) be a graded splitting of the anchor. Its failure to commute with the internal differential is
\[
 \tau_\sigma(\xi)
 :=d_{\mathfrak u}\sigma(\xi)-\sigma(d_{\bbT_S}\xi),
\]
and its bracket curvature is
\[
 F_\sigma(\xi,\eta)
 :=[\sigma(\xi),\sigma(\eta)]-\sigma([\xi,\eta]).
\]
These have bidegrees \((1,1)\) and \((2,0)\), respectively, and combine into the total curvature
\[
 \mathcal F_\sigma:=\tau_\sigma+F_\sigma
 \in
 \Tot^2 C^\bullet_{\mathrm{CE}}
 \bigl(\bbT_S;\mathfrak K_{\cF/S}\bigr).
\]

\begin{proposition}\label{prop:AKS-bianchi}
The covariant total differential \(D_\sigma\) determined by a graded splitting satisfies
\[
 D_\sigma^2=[\mathcal F_\sigma,-],
 \qquad
 D_\sigma\mathcal F_\sigma=0.
\]
If the vertical kernel is abelian, \(D_\sigma\) is a differential and \(\mathcal F_\sigma\) is a total cocycle.
\end{proposition}

\begin{proof}
The component of the first identity in internal degree is the compatibility of the differential with the bracket; the horizontal component is the ordinary curvature identity. The mixed component records the differential of the bracket curvature and the Chevalley--Eilenberg differential of \(\tau_\sigma\). Their sum is the graded Jacobi identity. Applying the same identity once more gives the total Bianchi equation. If the kernel is abelian, the adjoint action of the total curvature vanishes.
\end{proof}

\begin{corollary}\label{cor:AKS-obstructions}
The class \(\kappa_{\cF/S}\) is the obstruction to the existence of a derived transverse splitting of the anchor. Once a graded splitting exists, the obstruction to flatness is its total Maurer--Cartan curvature \(\mathcal F_\sigma\). In the abelian case the class
\[
 [\mathcal F_\sigma]
 \in
 H^2\Tot C^\bullet_{\mathrm{CE}}
 \bigl(\bbT_S;\mathfrak K_{\cF/S}\bigr)
\]
is independent of the graded splitting.
\end{corollary}

\begin{proof}
The first assertion is Proposition~\ref{prop:AKS-splitting}. If
\(\sigma'=\sigma+\alpha\), then
\[
 \mathcal F_{\sigma'}
 =\mathcal F_\sigma+D_\sigma\alpha
  +\frac12[\alpha,\alpha].
\]
When the kernel is abelian the final term vanishes and the differential is independent of the splitting, so the cohomology class is well-defined.
\end{proof}

\section{Existence of flat transverse splittings}\label{sec:existence}

Return to the affine effective controller
\[
 a_{\uuder}\from\uuder(\g/B)\longrightarrow T_{B/\kk}
\]
and put \(\kfr=\ker(a_{\uuder})\). Suppose the anchor is surjective. Since \(T_{B/\kk}\) is projective for \(B\) smooth, it admits a graded \(B\)-linear section.

\begin{lemma}\label{lem:graded-splitting-projective}
A surjection of graded \(B\)-modules onto a degreewise projective graded module admits a graded section. In particular, the preceding anchor admits a graded section.
\end{lemma}

\begin{proof}
Choose a section degree by degree. Smoothness of \(B\) makes \(\Omega^1_{B/\kk}\), and hence its dual \(T_{B/\kk}\), projective.
\end{proof}

For a graded section \(\sigma\), define \(\tau_\sigma\), \(F_\sigma\) and
\(\mathcal F_\sigma\) as above. If \(\alpha\) has total degree one, then
\[
 \mathcal F_{\sigma+\alpha}
 =\mathcal F_\sigma+D_\sigma\alpha
  +\frac12[\alpha,\alpha].
\]
For a strict correction of a section, \(\alpha\) lies in
\(\Omega^1_{B/\kk}\otimes_B\kfr^0\); the full derived problem permits all total degree-one components.

\begin{theorem}\label{thm:abelian-obstruction}
Assume that \(\kfr\) is abelian. Then
\[
 \mathfrak o(\g/B)
 :=[\mathcal F_\sigma]
 \in
 H^2\Tot C^\bullet_{\mathrm{CE}}
 \bigl(T_{B/\kk};\kfr\bigr)
\]
is independent of the graded section \(\sigma\). A derived flat splitting exists if and only if
\(\mathfrak o(\g/B)=0\). When it exists, the space of its equivalence classes is a torsor under
\[
 H^1\Tot C^\bullet_{\mathrm{CE}}
 \bigl(T_{B/\kk};\kfr\bigr).
\]
If \(\kfr\) is concentrated in degree zero, or more generally if every relevant total degree-one class has a representative of bidegree \((1,0)\), the same criterion classifies strict flat dg-Lie-algebroid splittings.
\end{theorem}

\begin{proof}
By Proposition~\ref{prop:AKS-bianchi}, \(\mathcal F_\sigma\) is a cocycle. Since \(\kfr\) is abelian, changing the section by \(\alpha\) changes the curvature by the coboundary \(D\alpha\). Thus the class is independent of \(\sigma\). It vanishes precisely when a total degree-one correction solves
\(\mathcal F_\sigma+D\alpha=0\), which is the abelian Maurer--Cartan equation. Two solutions differ by a closed total degree-one cochain, and equivalence changes this difference by an exact cochain. The final assertion identifies total solutions with strict corrections under the stated bidegree condition.
\end{proof}

\begin{remark}
For non-abelian kernels the existence problem is the curved Maurer--Cartan equation
\[
 \mathcal F_\sigma+D_\sigma\alpha
 +\frac12[\alpha,\alpha]=0.
\]
For the crossed controller it is the corresponding equation in the twisted total convolution \(L_\infty\)-algebra. Thus no reduction to an underived or effective deformation problem is required.
\end{remark}

\section{A split normal form for the unfolded de Rham algebra}\label{sec:normal-form}

Choose a graded splitting \(\sigma\) of the effective anchor. Write
\(R=\DR(\cF/S)\) and let \(d_{\mathrm{int}}\) and \(\epsilon_{\cF/S}\) be its internal and mixed differentials. A correction
\(\theta\in\Omega^1_{B/\kk}\otimes_B\kfr^0\) defines the transverse lift
\(\sigma+\theta\). On the Tate realisation of
\(R\otimes_B\DR(B/\kk)\), set
\[
 D_\theta
 :=d_{\mathrm{int}}
   +\epsilon_{\cF/S}+\epsilon_B
   +L_{\sigma+\theta},
\]
where the weight shifts implicit in the Tate realisation make the displayed sum homogeneous.

\begin{proposition}\label{prop:normal-form}
The internal and horizontal components of the square of \(D_\theta\) combine into
\[
 D_\theta^2
 =L_{\mathcal F_\sigma+D_\sigma\theta+\frac12[\theta,\theta]}.
\]
Equivalently, before Tate realisation one has the two graded-mixed identities
\[
 [d_{\mathrm{int}},\epsilon_\theta]
 =L_{\tau_{\sigma+\theta}},
 \qquad
 \epsilon_\theta^2=L_{F_{\sigma+\theta}},
\]
where
\(\epsilon_\theta=\epsilon_{\cF/S}+\epsilon_B+L_{\sigma+\theta}\).
If the vertical action is faithful, the resulting graded-mixed structure is integrable if and only if the total Maurer--Cartan curvature vanishes. Hence effective flat splittings are equivalent to split normal forms satisfying the total Maurer--Cartan equation.
\end{proposition}

\begin{proof}
The commutator with the internal differential measures the defect of \(\sigma+\theta\) as a chain map and gives the Lie derivative of \(\tau_{\sigma+\theta}\). The square of the mixed part measures failure to preserve the base bracket and gives the Lie derivative of \(F_{\sigma+\theta}\). The cross terms give the covariant differential of \(\theta\), while the quadratic term is \(\frac12[\theta,\theta]\). Summing the two bidegrees on the Tate realisation yields the displayed total formula. Faithfulness identifies vanishing of the Lie derivative with vanishing of its coefficient.
\end{proof}

\begin{proposition}\label{prop:higher-normal-form}
For the non-effective crossed controller, a split normal form is governed by a Maurer--Cartan element of the total convolution \(L_\infty\)-algebra associated with
\(\Uder(\g/B)\). Its unary component is the transverse Lie-derivative term above, while its higher components provide the coherent null-homotopies of central curvature. The full \(L_\infty\) Maurer--Cartan equation is equivalent to the square-zero condition for the resulting graded-mixed structure together with its prescribed coherences.
\end{proposition}

\begin{proof}
A representation of the two-term \(L_\infty\)-controller on the relative de Rham algebra sends the brackets of the convolution algebra to commutators of derivations and its higher brackets to the corresponding coherent homotopies. Twisting this representation by a Maurer--Cartan element changes the total differential by the associated unary and higher operators. The standard twisting identity identifies the square and all higher coherence defects with the image of the Maurer--Cartan curvature. Thus the twisted structure is integrable precisely when the full Maurer--Cartan equation holds \cite{Hinich01,Pridham10,Nuiten}.
\end{proof}

\begin{remark}
In the effective case the higher components vanish and Proposition~\ref{prop:higher-normal-form} reduces to Proposition~\ref{prop:normal-form}. In the presence of central isotropy, the equality \(D_\theta^2=0\) alone would only say that the curvature acts trivially; the higher component records its chosen null-homotopy and is therefore indispensable.
\end{remark}


\section{Leaf-space interpretation in the representable case}\label{sec:leaf-space}

The preceding constructions are internal to the foliation. When a leaf space is represented by a smooth morphism of schemes, the same theory admits a classical geometric interpretation: a flat transverse unfolding is an Ehresmann connection on the leaf space.

\begin{hypothesis}
Let
\[
\begin{tikzcd}
X \arrow[r,"q"] \arrow[dr,"\pi"'] & Y \arrow[d,"p"] \\
& S
\end{tikzcd}
\]
be smooth morphisms of smooth schemes, with \(q\) faithfully flat and \(\pi=p\circ q\). Let \(\cF/S\) be the relative foliation on \(X/S\) whose tangent complex is \(T_{X/Y}\). Equivalently, \(\cF/S\) is induced by the smooth leaf-space map \(q\).
\end{hypothesis}


\begin{proposition}\label{prop:leaf-controller}
Under the hypothesis above there are two distinct transverse objects.
First, the normal sheaf of the foliation induced by \(q\) is the \(\cO_X\)-module
\[
 T_X/T_{X/Y}\simeq q^*T_{Y/\kk}.
\]
Second, the sheaf of vector fields projectable along \(q\) is defined tautologically by
\[
 T_X^{\mathrm{proj}/Y}:=
 T_X\times_{\underline{\Der}_{\kk}(q^{-1}\cO_Y,\cO_X)}q^{-1}T_{Y/\kk}.
\]
Here the first arrow is restriction of derivations to \(q^{-1}\cO_Y\). The second arrow is obtained from a derivation of \(\cO_Y\) by applying \(q^{-1}\) and composing with \(q^{-1}\cO_Y\to\cO_X\). There is an exact sequence of sheaves of \(\kk\)-Lie algebras
\[
 0\longrightarrow T_{X/Y}\longrightarrow T_X^{\mathrm{proj}/Y}
 \longrightarrow q^{-1}T_{Y/\kk}\longrightarrow 0.
\]
Thus the projectable transverse symmetry sheaf is \(q^{-1}T_{Y/\kk}\), not \(q^*T_{Y/\kk}\). The latter is obtained from the former only after extension of scalars from \(q^{-1}\cO_Y\) to \(\cO_X\), and this extension forgets the projectability condition.

The transverse anchor for projectable symmetries is the inverse image of
\[
 dp\from T_{Y/\kk}\longrightarrow p^*T_{S/\kk}.
\]
Consequently, leaf-space-basic transverse splittings are those \(\cO_X\)-linear splittings of the normal sequence whose transverse classes descend from \(Y\); equivalently, they are the pullbacks along \(q\) of splittings of the tangent sequence of \(p\from Y\to S\).
\end{proposition}

\begin{proof}
The first assertion is the usual exact sequence of tangent sheaves for the smooth morphism \(q\):
\[
 0\longrightarrow T_{X/Y}\longrightarrow T_X
 \longrightarrow q^*T_{Y/\kk}\longrightarrow 0.
\]
It identifies the \(\cO_X\)-linear normal sheaf of the foliation with \(q^*T_{Y/\kk}\). This statement concerns normal directions as an \(\cO_X\)-module; it says nothing about whether a normal vector field is represented by a vector field on \(X\) which is related to a vector field on \(Y\).

By definition, a section of \(T_X^{\mathrm{proj}/Y}\) is a pair \((V,W)\), where \(V\) is a vector field on \(X\), \(W\) is a local vector field on \(Y\), and the derivation induced by \(V\) on functions pulled back from \(Y\) is the pullback of \(W\). Projection to the second component gives a morphism of sheaves of \(\kk\)-Lie algebras
\[
 T_X^{\mathrm{proj}/Y}\longrightarrow q^{-1}T_{Y/\kk}.
\]
Its kernel consists of those pairs \((V,0)\) for which \(V\) annihilates \(q^{-1}\cO_Y\); this is \(T_{X/Y}\). Since \(q\) is smooth, the morphism is locally surjective: after replacing \(X\) and \(Y\) by smooth local charts for \(q\), one may write \(X\simeq Y\times\mathbb A^r\), and a vector field \(W\) on \(Y\) is lifted to \((W,0)\) on \(Y\times\mathbb A^r\). Hence
\[
 0\to T_{X/Y}\to T_X^{\mathrm{proj}/Y}\to q^{-1}T_{Y/\kk}\to0
\]
is exact as a sequence of sheaves of \(\kk\)-Lie algebras.
The tautological definition is necessary. The condition \([V,T_{X/Y}]\subset T_{X/Y}\) is satisfied by every vector field when \(q\) is \'{e}tale, and therefore cannot by itself characterise vector fields related to vector fields on \(Y\). For example, for \(q\from\mathbb G_m\to\mathbb G_m\), \(y=x^2\), the relative tangent sheaf is zero, but \(x^2\partial_x\) is not related to a vector field on the target, since its action on \(y\) is \(2x^3\), not the pullback of a regular function of \(y\).
The distinction between \(q^{-1}T_{Y/\kk}\) and \(q^*T_{Y/\kk}\) is also visible in the elementary projection \(\mathbb A^2_{x,y}\to\mathbb A^1_y\). The normal module contains \(x\partial_y\), but no vector field representing this normal class is projectable, since
\[
 [x\partial_y+g(x,y)\partial_x,\partial_x]
 =-\partial_y-\frac{\partial g}{\partial x}\partial_x
\]
is not vertical. Thus projectability is not \(\cO_X\)-linear.
The map on projectable transverse symmetries is the inverse image of \(dp\). Extending scalars gives the normal morphism
\[
 q^*T_{Y/\kk}\longrightarrow q^*p^*T_{S/\kk}=\pi^*T_{S/\kk}.
\]
A transverse normal splitting is called leaf-space-basic when this \(\cO_X\)-linear splitting is obtained by pulling back a splitting of
\[
 T_{Y/\kk}\longrightarrow p^*T_{S/\kk}.
\]
Faithfully flat descent for quasi-coherent sheaves along \(q\) identifies such descended splittings with splittings on \(Y\), as claimed.
\end{proof}

\begin{theorem}\label{thm:leaf-space-connection}
In the representable leaf-space situation, leaf-space-basic transversal unfoldings of \(\cF/S\) are equivalent to flat splittings
\[
 p^*T_{S/\kk}\longrightarrow T_{Y/\kk}
\]
of the tangent sequence of \(p\from Y\to S\). Equivalently, they are integrable Ehresmann connections on the leaf-space fibration \(Y\to S\), pulled back to \(X\).
\end{theorem}

\begin{proof}
By Proposition \ref{prop:leaf-controller}, the \(\cO_X\)-linear normal sheaf is \(q^*T_{Y/\kk}\), whereas the projectable transverse symmetries before extension of scalars form \(q^{-1}T_{Y/\kk}\). Hence a leaf-space-basic splitting is not an arbitrary splitting of
\[
 q^*T_{Y/\kk}\longrightarrow q^*p^*T_{S/\kk},
\]
but a splitting carrying descent data along \(q\), equivalently the scalar extension of a unique splitting of
\[
 T_{Y/\kk}\longrightarrow p^*T_{S/\kk}.
\]
This is an Ehresmann connection for \(p\from Y\to S\).

The bracket curvature of the pulled-back splitting is the pullback of the bracket curvature of the splitting on \(Y\). Since \(q\) is faithfully flat, the pulled-back curvature vanishes if and only if the curvature on \(Y\) vanishes. Therefore flat leaf-space-basic transverse splittings are flat Ehresmann connections on \(Y\to S\). Applying the affine classification theorem to these descended projectable splittings gives the announced equivalence with leaf-space-basic transversal unfoldings.
\end{proof}

\begin{remark}
Theorem \ref{thm:leaf-space-connection} explains the role of the analytic leaf spaces appearing in the theory of derived foliations. Whenever a derived foliation admits an analytic leaf space satisfying the hypotheses of Section~\ref{sec:leaf-space}, a leaf-space-basic flat transverse unfolding is viewed here as a flat connection on that leaf space. The intrinsic controller \(\Uder(\cF/S)\) is the replacement for \(T_Y\) when the leaf space is formal, derived, singular, or not represented by a scheme.
\end{remark}



\section{Filtered Gauss--Manin connection and characteristic classes}\label{sec:filtered-GM}

In the present Cartan-linearised setting the transverse operators have weight zero. Consequently the Gauss--Manin connection constructed above preserves the weight filtration on the Tate realisation. We shall refer to the resulting filtration as the Cartan-linearised Hodge filtration. Its preservation is stronger than the usual Griffiths-transversality statement for classical Gauss--Manin connections and depends upon the chosen transverse Cartan linearisation; no corresponding assertion is intended for arbitrary filtered de Rham complexes. We record the argument because this is the form in which the construction interacts with foliated characteristic classes.

\begin{definition}
Let \(M\) be a coherently \(s\)-Cartan-linearised crystal along \(\cF/S\). The Hodge filtration on
\[
 \DR_{\cF/S}(M)=M^{\Tate}
\]
is the complete decreasing filtration induced by the weight filtration of the graded mixed module \(M\). It is denoted
\[
 F^p\DR_{\cF/S}(M).
\]
After applying \(R\pi_*\), we set
\[
 F^pR\pi_*\DR_{\cF/S}(M):=R\pi_*F^p\DR_{\cF/S}(M).
\]
\end{definition}

\begin{proposition}\label{prop:filtered-GM}
Under the hypotheses of Theorem~\ref{thm:GM}, and for the Cartan-linearised weight/Hodge filtration just defined, the Gauss--Manin connection preserves the filtration:
\[
 \nabla^{\GM}_s
 \bigl(F^pR\pi_*\DR_{\cF/S}(M)\bigr)
 \subset
 \bbL_S\otimes F^pR\pi_*\DR_{\cF/S}(M).
\]
Thus \(R\pi_*\DR_{\cF/S}(M)\) is a filtered crystal over \(S\).
\end{proposition}

\begin{proof}
We verify the assertion before descent. In a CE chart, the connection is induced by operators
\[
 L^M_{\Theta_\xi}\from M\to M
\]
where \(\Theta_\xi\) is a weight-zero graded mixed derivation of \(\DR(\cF/S)\). Since \(\Theta_\xi\) has weight zero, it maps the weight piece \(M^{(r)}\) into \(M^{(r)}\). Therefore, for the decreasing filtration
\[
 F^pM^{\Tate}=\prod_{r\ge p}M^{(r)}[-r]
\]
one has
\[
 L^M_{\Theta_\xi}(F^pM^{\Tate})\subset F^pM^{\Tate}.
\]
The mixed differential of \(M\) and the lifted Lie derivative commute because the linearisation is a graded mixed module action. Hence the inclusion is an inclusion of filtered complexes, not only of graded objects. Applying \(R\pi_*\) gives an endomorphism
\[
 R\pi_*(L^M_{\Theta_\xi})
 \from R\pi_*\DR_{\cF/S}(M)
 \longrightarrow R\pi_*\DR_{\cF/S}(M)
\]
whose restriction satisfies
\[
 R\pi_*(L^M_{\Theta_\xi})
 \bigl(F^pR\pi_*\DR_{\cF/S}(M)\bigr)
 \subset F^pR\pi_*\DR_{\cF/S}(M).
\]
The local operators glue by the descent argument in Theorem~\ref{thm:GM}; consequently the descended Gauss--Manin connection preserves the descended filtration.
\end{proof}

\begin{proposition}\label{prop:chern-horizontal}
Let \(E\) be a perfect coherently \(s\)-Cartan-linearised crystal along \(\cF/S\). Suppose that the foliated Chern character
\[
 \operatorname{ch}_{\cF/S}(E)
 \in
 H^0\bigl(S,R\pi_*\DR_{\cF/S}\bigr)
\]
is defined by a construction on perfect Cartan-linearised crystals which is functorial under equivalences, compatible with restriction to CE charts, and invariant under Cartan homotopies of tangent-inner actions. Then
\[
 \nabla^{\GM}_s\operatorname{ch}_{\cF/S}(E)=0.
\]
The same assertion holds for the foliated Chern classes obtained from the Chern character by the universal Newton identities.
\end{proposition}

\begin{proof}
The assertion follows from functoriality, rather than from a separate curvature computation. The flat splitting \(s\) gives, by Theorem~\ref{thm:GM}, parallel transport operators on the relative foliated de Rham complex of every coherently \(s\)-Cartan-linearised crystal. In a CE chart and for \(\xi\in T_{B/\kk}\), choose a representative \(\Theta_\xi\) of \(s(\xi)\). The \(s\)-linearisation gives a Lie derivative \(L^E_{\Theta_\xi}\) on \(E\), hence an infinitesimal automorphism of the corresponding perfect object in the Cartan-linearised crystal category.

By functoriality of the chosen Chern character construction, applying \(L^E_{\Theta_\xi}\) to \(E\) differentiates \(\operatorname{ch}_{\cF/S}(E)\) by the induced Lie derivative on the target foliated de Rham complex. An infinitesimal automorphism acts trivially on the functorial Chern character after passage to cohomology: in strict CE models this is the usual trace-of-a-commutator calculation, and in the abstract formulation it is the assumed functoriality and homotopy invariance of the Chern character. Therefore the covariant derivative of \(\operatorname{ch}_{\cF/S}(E)\) along \(\xi\) vanishes.

It remains to check independence of choices. If another representative of \(s(\xi)\) is chosen, the difference is tangent-inner. The Cartan-linearisation supplies the contraction homotopy of Lemma~\ref{lem:cartan-homotopy}; hence the two induced actions on foliated de Rham cohomology agree. The compatibility of the Chern character with CE restrictions makes the local argument descend along the chosen hypercover. Thus
\[
 \nabla^{\GM}_{s,\xi}\operatorname{ch}_{\cF/S}(E)=0
\]
for every local vector field \(\xi\), and the global horizontality follows. Since the ordinary Chern classes are universal polynomials in the components of the Chern character via the Newton identities, their horizontality follows formally.
\end{proof}

\begin{corollary}
Assume \(\pi\) is proper and that the trace map on relative foliated de Rham cohomology is defined and compatible with the Gauss--Manin connection. Numerical invariants obtained by integrating foliated characteristic classes of perfect Cartan-linearised crystals along the fibres are locally constant under a flat transverse unfolding.
\end{corollary}

\begin{proof}
By Proposition~\ref{prop:chern-horizontal}, the relevant characteristic classes are horizontal. Applying the trace map gives horizontal sections of \(\cO_S\). A horizontal section for the ordinary de Rham connection is locally constant on each connected component of \(S\).
\end{proof}

\section{Examples and illustrations}\label{sec:examples}

The examples of this section illustrate the geometric content of the classification theorem in several familiar settings. In the classical and logarithmic cases the controller is a familiar quotient of projectable vector fields; in the representable leaf-space case it is the tangent algebroid of the leaf-space fibration; in derived pullback and formal-leaf situations it is the homotopical replacement for that tangent object; and for Cartan-linearised crystals it is the source of Gauss--Manin transport. Thus the same flat transverse splitting governs the examples from classical transverse symmetries to filtered foliated cohomology.

\subsection{Classical smooth foliations}

Let \(X\to S\) be smooth and let \(T_{\cF/S}\subset T_{X/S}\) be a subbundle closed under the Lie bracket. The associated relative derived foliation is underived: its de Rham algebra is
\[
 \DR(\cF/S)=\Sym_{\cO_X}(\Omega^1_{\cF/S}[1]),
\]
with mixed differential equal to the ordinary leafwise de Rham differential. Let
\[
 q_\pi\from T_X\longrightarrow \pi^*T_{S/\kk}
\]
be the quotient map. The effective controller becomes
\[
 \uuder(\cF/S)=
 \frac{
 \{Y\in T_X:
 q_\pi(Y)\text{ lies in the image of }\pi^{-1}T_{S/\kk}\to\pi^*T_{S/\kk},
 [Y,T_{\cF/S}]\subset T_{\cF/S}\}
 }{T_{\cF/S}} .
\]
Thus a flat splitting is a projectable flat system of vector fields preserving the foliation, modulo fields tangent to the leaves. This recovers the classical transverse-symmetry interpretation of an unfolding.

\subsection{Tautological derived foliations}

Let \(X=\Spec A\). To\"en--Vezzosi distinguish two tautological objects: the pointwise foliation and the de Rham foliation. The pointwise foliation is obtained by the augmentation \(\DR(A/\kk)\to A\); the de Rham foliation is obtained from the identity of \(\DR(A/\kk)\). They are respectively the initial and final derived foliations on \(X\).

They are not direct applications of the strict theorem. Following the tautological examples in \cite[Section~2.1]{TV}, they serve as calibration cases for the definition of the controller. In the pointwise foliation there is no leafwise direction and every ambient direction is transverse. In the de Rham foliation every ambient direction is leafwise and there is no transverse quotient. Any intrinsic definition of \(\Uder(\cF/S)\) must reproduce these two limiting cases.

\subsection{Relative de Rham foliations induced by a morphism}

Let \(f\from X=\Spec A\to Y=\Spec B\) be a morphism of derived affine schemes. The relative de Rham foliation is defined by
\[
 \DR(A/B)=\DR(A/\kk)\otimes_{\DR(B/\kk)}B.
\]
Its cotangent complex is \(\bbL_{A/B}\). The cotangent triangle
\[
 A\otimes_B\bbL_{B/\kk}\longrightarrow \bbL_{A/\kk}\longrightarrow \bbL_{A/B}
\]
is the most elementary incarnation of the distinction between transverse and relative directions.

When \(A\) and \(B\) are smooth and \(f\) is a submersion, the foliation is the ordinary foliation by the fibres. A flat transversal unfolding is then a flat Ehresmann connection on the family. If \(f\) is not smooth, the same construction still defines a derived foliation; the higher homology of \(\bbL_{A/B}\) measures the failure of the classical tangent distribution to capture the fibre geometry.

\subsection{Ordinary Lie algebroids}

Let \(\rho\from\g\to T_{A/B}\) be an ordinary vector-bundle Lie algebroid in the Lie--Rinehart sense \cite{Rinehart63}. Its Chevalley--Eilenberg algebra
\[
 \CE^*(\g)=\Sym_A(\g^\vee[1])
\]
with the CE differential is a derived foliation. If \(\rho\) is injective and identifies \(\g\) with a subbundle of \(T_{A/B}\), the preceding smooth foliation case is recovered. If \(\rho\) has isotropy, the derived foliation retains it through the CE algebra. The crossed-module controller \([\g\to D^1_{\bas}(\g/B)]\) is then not an auxiliary refinement, but the minimal structure which retains the stabilisers of unfoldings.
This is the derived extension of the earlier theorem on unfoldings of singular holomorphic Lie algebroids \cite{CMQ}: the classical controller is recovered after passage to the effective underived truncation.

\subsection{The Suwa transverse term}

Let \(S=\Spec \kk[\eps]/(\eps^2)\). In codimension one, Suwa's local expression for an unfolding \cite{Suwa81,Suwa83} is
\[
 \omega+\eps\eta+h\,d\eps.
\]
The term \(h\,d\eps\) is the lift of the infinitesimal parameter direction. In the present notation it is the value of a transverse splitting on \(\partial_\eps\):
\[
 s(\partial_\eps)\in \uuder(\g/B).
\]
The integrability of the unfolded form is the flatness of this splitting. In the split graded-mixed normal form this becomes
\[
 \epsilon_\theta=\epsilon_{\cF/S}+\epsilon_B+L_\theta,
 \qquad
 \epsilon_\theta^2=0.
\]
Thus Suwa's formula is the one-dimensional classical trace of the Maurer--Cartan equation for the unfolded mixed differential.

\subsection{Derived pullbacks of smooth foliations}

Let \(\cG\) be a smooth classical foliation on a smooth scheme \(Y\), and let \(f\from X\to Y\) be any morphism. The functoriality of derived foliations \cite[Section~2.1]{TV} gives a pullback derived foliation \(f^*\cG\) on \(X\). If \(f\) is transverse to \(\cG\), the pullback is classical. If transversality fails, the pullback generally has a non-trivial cotangent complex and cannot be represented by a vector subbundle of \(T_X\).

This example shows that a controller defined solely from a singular classical distribution is insufficient. The intrinsic datum is the graded-mixed algebra \(\DR(f^*\cG)\), and the transverse controller must be computed from its basic graded-mixed derivations.




\subsection{Shifted Poisson structures and derived symplectic foliations}

The shifted Poisson example must be read in the same homotopical sense as the rest of the paper. A shifted Poisson structure is naturally encoded by a Hamiltonian \(L_\infty\)-algebroid, or equivalently by a suitable \(P_{n+1}\)-algebraic structure, rather than by a strict dg-Lie algebroid without choices. Assume therefore that the given shifted Poisson structure on \(X\) is represented by a strict Hamiltonian \(L_\infty\)- or CE presentation \(\g_\pi\), whose underlying complex is the shifted cotangent complex and whose anchor is induced by the Poisson tensor. The associated Chevalley--Eilenberg graded mixed algebra then defines the derived symplectic foliation to which the preceding theory applies.

For a family \(X\to S\), a transversal unfolding of this Hamiltonian presentation is a flat deformation of the derived symplectic foliation along the base. Its transverse controller is the basic-derivation controller of \(\g_\pi\), interpreted as a crossed controller when the Hamiltonian presentation has isotropy. Thus the classification theorem applies to strict Hamiltonian presentations and is invariant under replacement by an equivalent presentation through the intrinsic graded-mixed controller.

In the non-degenerate shifted symplectic case, the Hamiltonian foliation is the de Rham foliation. The controller then reduces to the customary algebroid governing flat Ehresmann transport in the family. In degenerate shifted Poisson situations the same formula records the transverse symmetries of the derived symplectic leaves, including higher stabilisers coming from the Hamiltonian \(L_\infty\)-structure.

\subsection{Quasi-smooth singular enhancements}

A quasi-smooth rigid derived foliation on a smooth scheme, in the sense used by To\"en--Vezzosi \cite{TVRH,TV}, has a cotangent complex sitting locally in a triangle
\[
 N^*_{\cF}\longrightarrow\Omega^1_X\longrightarrow\bbL_\cF,
\]
with \(N^*_{\cF}\) a vector bundle. On the open locus where \(\bbL_\cF\) is a vector bundle, one obtains an ordinary smooth foliation; the complement is the singular locus. The derived enhancement records the homotopical correction along this singular locus.

The unfolding theorem applies after choosing local CE presentations satisfying the hypotheses of the paper. Even before such presentations are chosen, the example explains the intrinsic nature of the controller: the singular truncation alone does not contain the coherent mixed differential, and therefore cannot determine the correct transverse symmetries.

\subsection{Formal leaves and formal leaf spaces}

Almost perfect derived foliations on derived Artin stacks locally of finite presentation are formally integrable around points \cite[Section~2.3]{TV}. Hence, in that setting, a point carries a formal leaf and a formal leaf space. This statement is not available for arbitrary classical singular distributions: it is an intrinsically derived integrability statement, and it is one reason why not every singular classical foliation should be expected to admit a derived enhancement.

This leads to the following interpretation: a transversal unfolding is a flat infinitesimal transport of the formal leaf space along the base. If the formal leaf space is not represented by a scheme or analytic space, the controller \(\Uder(\cF/S)\) is the infinitesimal replacement for the tangent object of that missing leaf space.

\subsection{Holomorphic analytic leaf spaces}

In the holomorphic theory, quasi-smooth rigid derived foliations are formally integrable and, under suitable analytic hypotheses, locally analytically integrable \cite[Chapters~5--6]{TV}. To\"en--Vezzosi use this to construct holomorphic leaf spaces in favourable cases \cite[Section~6.2]{TV}. One obtains a flat surjective holomorphic map
\[
 q\from X\longrightarrow X/\cF
\]
to a possibly highly non-separated complex orbifold, and \(\cF\) is recovered as the derived foliation induced by \(q\).

If a relative holomorphic derived foliation \(\cF/S\) admits such a relative leaf space
\[
 X\xto{q}Y\xto{p}S,
\]
then the representable leaf-space result of Section~\ref{sec:leaf-space} applies to leaf-space-basic unfoldings: a flat transversal unfolding is a flat Ehresmann connection on \(Y\to S\), pulled back to \(X\). When the leaf space is non-separated or only formal, the same interpretation remains valid only after replacing the tangent sheaf \(T_Y\) by the intrinsic homotopical controller \(\Uder(\cF/S)\).

\subsection{Cartan-linearised crystals and filtered Gauss--Manin cohomology}

A quasi-coherent crystal along \(\cF/S\) is a graded-mixed module induced from its quasi-coherent weight-zero component as in Definition~\ref{def:crystals}, and its Tate realisation is its foliated de Rham complex \cite[Sections~3.1--3.2]{TV}. A transverse unfolding acts by Gauss--Manin transport on crystals which are linearised by the unfolding and Cartan-linearised in the sense of Definition~\ref{def:linearised-crystal}. For such crystals the Gauss--Manin connection of Section~\ref{sec:crystals} acts on relative foliated de Rham cohomology. Since the transverse Cartan operators have weight zero, it preserves the Cartan-linearised weight/Hodge filtration defined in Section~\ref{sec:filtered-GM}.

Thus the construction is a consequence of the unfolding theory established here. The unfolded mixed differential supplies transverse Lie derivatives on the coefficient object; the Cartan homotopy kills changes by tangent inner derivations; flatness gives the curvature-zero condition on the resulting connection. Accordingly, Cartan-linearised crystals are transported along the base by a flat transversal unfolding.

\subsection{A non-effective example with central isotropy}

Let \(\g=\h\oplus\z\), where \(\h\to T_{A/B}\) is an effective relative dg-Lie algebroid and \(\z\) is a complex of central abelian Lie algebras with zero anchor. Then
\[
 \ker(\iota_\g)=\z.
\]
The ordinary quotient \(D^1_{\bas}(\g/B)/\iota(\g)\) forgets automorphisms coming from \(\z\). The crossed-module controller
\[
 [\h\oplus\z\to D^1_{\bas}(\g/B)]
\]
retains \(\z\) as pointwise isotropy. For a fixed splitting, the strict automorphisms are the Chevalley--Eilenberg cocycles with values in the transported module \(\z_s\); the full derived stabiliser is governed by the total complex \(\Tot C^\bullet_{\mathrm{CE}}(T_{B/\kk};\z_s)\). These automorphisms and higher homotopies disappear from the ordinary effective quotient but remain visible in the crossed classification.

\subsection{Logarithmic tangent algebroids}

Let \((X,D)\to S\) be log-smooth in the standard logarithmic-geometric sense of Kato \cite{Kato89}, with \(D\subset X\) a relative normal crossing divisor; if horizontal components are present, replace \(T_{S/\kk}\) below by the corresponding logarithmic tangent sheaf of the base. The logarithmic relative tangent sheaf
\[
 T_{X/S}^{\log}=\operatorname{Der}_{\cO_S}(\cO_X,\cO_X)(-\log D)
\]
is a relative Lie algebroid. A transversal unfolding is a flat projectable system of logarithmic vector fields preserving the relative logarithmic algebroid. Equivalently, it is a flat logarithmic Ehresmann connection compatible with the log structure. The controller is the quotient
\[
\frac{
\left\{
\begin{array}{c|l}
Y\in T_X^{\log} &
\begin{array}{l}
q_\pi(Y)\in \operatorname{im}\bigl(\pi^{-1}T_{S/\kk}\to\pi^*T_{S/\kk}\bigr),\\[2pt]
[Y,T_{X/S}^{\log}]\subset T_{X/S}^{\log}
\end{array}
\end{array}
\right\}}
{T_{X/S}^{\log}}.
\]
The resulting Gauss--Manin connection acts on logarithmic foliated de Rham cohomology; in the Cartan-linearised case its transverse operators have weight zero and preserve the logarithmic weight/Hodge filtration.



\begin{thebibliography}{99}

\bibitem{Atiyah57}
M. F. Atiyah,
\emph{Complex analytic connections in fibre bundles},
Trans. Amer. Math. Soc. \textbf{85} (1957), 181--207.

\bibitem{CPTVV}
D. Calaque, T. Pantev, B. To\"en, M. Vaqui\'e and G. Vezzosi,
\emph{Shifted Poisson structures and deformation quantization},
J. Topol. \textbf{10} (2017), no. 2, 483--584.


\bibitem{BrownSpencer76}
R. Brown and C. B. Spencer,
\emph{G-groupoids, crossed modules and the fundamental groupoid of a topological group},
Indag. Math. \textbf{38} (1976), no. 4, 296--302.

\bibitem{Deligne71}
P. Deligne,
\emph{Th\'eorie de Hodge. II},
Inst. Hautes \`Etudes Sci. Publ. Math. \textbf{40} (1971), 5--57.

\bibitem{Deligne74}
P. Deligne,
\emph{Th\'eorie de Hodge. III},
Inst. Hautes \`Etudes Sci. Publ. Math. \textbf{44} (1974), 5--77.

\bibitem{CMQ}
M. Corr\^ea, A. Molinuevo and F. Quallbrunn,
\emph{Families and unfoldings of singular holomorphic Lie algebroids},
Ann. Sc. Norm. Super. Pisa Cl. Sci. (5) \textbf{24} (2023), no. 2, 989--1006.

\bibitem{Hinich01}
V. Hinich,
\emph{DG coalgebras as formal stacks},
J. Pure Appl. Algebra \textbf{162} (2001), no. 2--3, 209--250.

\bibitem{HAGII}
B. To\"en and G. Vezzosi,
\emph{Homotopical algebraic geometry. II. Geometric stacks and applications},
Mem. Amer. Math. Soc. \textbf{193} (2008), no. 902.

\bibitem{IllusieI}
L. Illusie,
\emph{Complexe cotangent et d\'eformations. I},
Lecture Notes in Mathematics, Vol. 239, Springer-Verlag, Berlin--New York, 1971.

\bibitem{KatzOda68}
N. M. Katz and T. Oda,
\emph{On the differentiation of de Rham cohomology classes with respect to parameters},
J. Math. Kyoto Univ. \textbf{8} (1968), 199--213.

\bibitem{Kato89}
K. Kato,
\emph{Logarithmic structures of Fontaine--Illusie},
in \emph{Algebraic Analysis, Geometry, and Number Theory},
Johns Hopkins Univ. Press, Baltimore, 1989, pp. 191--224.

\bibitem{Malgrange76}
B. Malgrange,
\emph{Frobenius avec singularit\'es. I. Codimension un},
Inst. Hautes \`Etudes Sci. Publ. Math. \textbf{46} (1976), 163--173.

\bibitem{Malgrange77}
B. Malgrange,
\emph{Frobenius avec singularit\'es. II. Le cas g\'en\'eral},
Invent. Math. \textbf{39} (1977), no. 1, 67--89.

\bibitem{Nuiten}
J. Nuiten,
\emph{Homotopical algebra for Lie algebroids},
Appl. Categ. Structures \textbf{27} (2019), no. 5, 493--534.

\bibitem{NuitenKoszul}
J. Nuiten,
\emph{Koszul duality for Lie algebroids},
arXiv:1712.03442.

\bibitem{Noohi07}
B. Noohi,
\emph{Notes on 2-groupoids, 2-groups and crossed modules},
Homology Homotopy Appl. \textbf{9} (2007), no. 1, 75--106.

\bibitem{PTVV}
T. Pantev, B. To\"en, M. Vaqui\'e and G. Vezzosi,
\emph{Shifted symplectic structures},
Publ. Math. Inst. Hautes \`Etudes Sci. \textbf{117} (2013), 271--328.

\bibitem{Pridham10}
J. P. Pridham,
\emph{Unifying derived deformation theories},
Adv. Math. \textbf{224} (2010), no. 3, 772--826.

\bibitem{Rinehart63}
G. S. Rinehart,
\emph{Differential forms on general commutative algebras},
Trans. Amer. Math. Soc. \textbf{108} (1963), 195--222.

\bibitem{Quallbrunn}
F. Quallbrunn,
\emph{Isotrivial unfoldings and structural theorems for foliations on projective spaces},
Bull. Braz. Math. Soc. (N.S.) \textbf{48} (2017), no. 3, 335--345.

\bibitem{Suwa81}
T. Suwa,
\emph{A theorem of versality for unfoldings of complex analytic foliation singularities},
Invent. Math. \textbf{65} (1981/82), no. 1, 29--48.

\bibitem{Suwa83}
T. Suwa,
\emph{Unfoldings of complex analytic foliations with singularities},
Japan. J. Math. (N.S.) \textbf{9} (1983), no. 1, 181--206.

\bibitem{TVRH}
B. To\"en and G. Vezzosi,
\emph{Algebraic foliations and derived geometry: the Riemann--Hilbert correspondence},
Selecta Math. (N.S.) \textbf{29} (2023), no. 1, Paper No. 5.

\bibitem{TVGRR}
B. To\"en and G. Vezzosi,
\emph{Algebraic foliations and derived geometry. II. The Grothendieck--Riemann--Roch theorem},
arXiv:2007.09251.

\bibitem{TV}
B. To\"en and G. Vezzosi,
\emph{Derived foliations},
arXiv:2305.08212v2, 2025.

\end{thebibliography}

\end{document}
